\chapter{Complexity — Cfg Record}
%%% generated by the blueprint pipeline from TCSlib.Complexity.TuringMachine.Build.CfgRecord
\section{Overview}

Declarations of the Lean module \texttt{TCSlib.Complexity.TuringMachine.Build.CfgRecord}.

\section{Declarations}

\begin{definition}[Configuration within a work-tape window]
\lean{Turing.Cfg.InWindow}
\label{Turing.Cfg.InWindow}
\leanok
\uses{Turing.Cfg}
Let a configuration of a $k$-tape Turing machine on input $x$ have work-head positions
$p_i$ and work-tape contents $T_i : \mathbb{Z}\to\mathrm{Option}(\mathrm{Symbol})$. For
a natural number $s$, the configuration is within the window of radius $s$ if every
work-head position satisfies $|p_i|\le s$, and every tape cell $z$ with $|z|>s$ is blank
on every work tape.
\end{definition}

\begin{definition}[Symbols for a machine record]
\lean{Turing.RecSym}
\label{Turing.RecSym}
\leanok
For types $Q$ and $T$, let the record alphabet consist of four kinds of symbols: a
binary digit carrying a Boolean value; a work-cell symbol carrying an optional Boolean
content and a Boolean indicating whether the head is there; a state symbol carrying an
optional state in $Q$; and a tag symbol carrying an element of $T$.
\end{definition}

\begin{definition}[Embedding of binary digits]
\lean{Turing.RecSym.digitEmb}
\label{Turing.RecSym.digitEmb}
\leanok
\uses{Turing.RecSym}
For types $Q$ and $T$, the Boolean values embed into the record alphabet as digit
symbols, with each Boolean sent to the digit symbol carrying that value.
\end{definition}

\begin{definition}[Binary digits of a natural number]
\lean{Turing.binDigits}
\label{Turing.binDigits}
\leanok
For natural numbers $w$ and $v$, the list of binary digits of $v$ of width $w$ is the
list, in increasing order of bit position from $0$ to $w-1$, whose entry at position $i$
is the $i$th binary digit of $v$.
\end{definition}

\begin{definition}[Width for encoding a position]
\lean{Turing.posWidth}
\label{Turing.posWidth}
\leanok
For a natural number $n$, define its position width to be the number of binary digits in
the base-two representation of $n+1$.
\end{definition}

\begin{definition}[Record encoding of a configuration]
\lean{Turing.cfgRecord}
\label{Turing.cfgRecord}
\leanok
\uses{Turing.Cfg, Turing.RecSym, Turing.binDigits, Turing.posWidth}
Let $k,s\text{ be natural numbers}$, let $Q,T$ be types, let $x$ be a finite Boolean
string, and let $c$ be a configuration of a $k$-tape machine on input $x$, with tag
$\tau\text{ in }T$. Its record is a list of fields: first the current state; then the
binary digits of the shifted input-head position, using width equal to the binary length
of $|x|+1$; then, for each work tape, the cells at positions $-s,-s+1,\\ldots,s$, each
recording its content and whether its head is there; and finally a field containing
$\tau$.
\end{definition}

\begin{theorem}[Length of a configuration record]
\lean{Turing.cfgRecord_length}
\label{Turing.cfgRecord_length}
\leanok
\uses{Turing.Cfg, Turing.cfgRecord}
\difficulty{1}
For any natural number $s$, configuration $c$ of a $k$-tape machine with Boolean
work-tape symbols, and tag $\tau$, the record encoding $c$ with tag $\tau$ has exactly
$k+3$ fields.
\end{theorem}

\begin{theorem}[Binary width bound]
\lean{Turing.record_pos_lt}
\label{Turing.record_pos_lt}
\leanok
\uses{Turing.posWidth}
\difficulty{3}
For every natural number $n$, let $w$ be the length of the binary expansion of $n+1$.
Then $n+1<2^w$.
\end{theorem}

\begin{theorem}[Value of a bounded binary digit list]
\lean{Turing.record_digits_value}
\label{Turing.record_digits_value}
\leanok
\uses{Turing.binDigits}
\difficulty{6}
Let $w$ and $v$ be natural numbers with $v<2^w$. The base-two value of the list of the
first $w$ binary digits of $v$ is $v$.
\end{theorem}

\begin{theorem}[Uniqueness of bounded binary digit lists]
\lean{Turing.record_digits_inj}
\label{Turing.record_digits_inj}
\leanok
\uses{Turing.binDigits, Turing.record_digits_value}
\difficulty{1}
Let $w,u,v$ be natural numbers with $u<2^w$ and $v<2^w$. If the lists of the first $w$
binary digits of $u$ and $v$ agree, then $u=v$.
\end{theorem}

\begin{theorem}[Injectivity of configuration records]
\lean{Turing.cfgRecord_inj}
\label{Turing.cfgRecord_inj}
\leanok
\uses{Turing.Cfg, Turing.Cfg.InWindow, Turing.binDigits, Turing.cfgRecord, Turing.posWidth, Turing.record_digits_inj, Turing.record_pos_lt}
\difficulty{5}
Fix a natural number $s$ and two configurations $c,d$ on the same finite Boolean input,
together with tags $\tau,\tau'$. Suppose the head positions of both configurations lie
within distance $s$ of the origin on every work tape, and both configurations have blank
cells at every integer position whose absolute value exceeds $s$. Their records with
tags $\tau$ and $\tau'$ are equal if and only if their states, input-head positions,
work-tape contents, and work-tape head positions agree, and their tags agree.
\end{theorem}

\begin{definition}[Validity of a configuration record]
\lean{Turing.RecordValid}
\label{Turing.RecordValid}
\leanok
\uses{Turing.Cfg, Turing.Cfg.InWindow, Turing.RecSym, Turing.cfgRecord}
Let $k,n,s$ be natural numbers, let $Q,T$ be types, and let $F$ be a finite list of
finite fields over the record alphabet for states $Q$ and tags $T$. The record $F$ is
valid if there exist a Boolean input $x$ of length $n$, a configuration $c$ on $x$ whose
work-tape heads lie in the window of radius $s$, and a tag $\tau$ in $T$ such that the
configuration record of $c$ with tag $\tau$ is $F$.
\end{definition}

\begin{theorem}[Characterization of valid configuration records]
\lean{Turing.recordValid_iff}
\label{Turing.recordValid_iff}
\leanok
\uses{Turing.Cfg, Turing.Cfg.InWindow, Turing.RecSym, Turing.RecordValid, Turing.binDigits, Turing.cfgRecord, Turing.cfgRecord_length, Turing.posWidth}
\difficulty{7}
For natural numbers $k,n,s$, types $Q,T$, and a list $f$ of lists of recording symbols,
$f$ is a valid record for some Boolean input of length $n$, a configuration whose
work-tape heads and nonblank cells lie within distance $s$ of the origin, and some tag
if and only if $f$ has length $k+3$; its first field is a singleton state symbol; its
second field is the fixed-width binary digit list of some natural number $v\le n+1$,
with width equal to the number of binary digits in $n+1$; each of the next $k$ fields
lists $2s+1$ tape cells, with arbitrary optional Boolean contents and a mark on exactly
one position $m<2s+1$; and its final field is a singleton tag symbol.
\end{theorem}

\begin{definition}[Work tape representation of a field record]
\lean{Turing.RecordAt}
\label{Turing.RecordAt}
\leanok
\uses{Turing.Cfg, Turing.fieldsTape, Turing.joinFields}
Let $d$ be a configuration of a $K$-tape Turing machine, let $r$ index one of its work
tapes, and let $F$ be a finite list of finite lists of tape symbols. We say that $d$
represents $F$ on tape $r$ if, for every integer $p$ with $-1 \le p \le
|\operatorname{join}(F)|$, the symbol or blank stored on tape $r$ at the position $p$
units from its head equals the tape content specified by $F$ at offset $p$.
\end{definition}

\begin{definition}[Bounded run between configurations]
\lean{Turing.RecordRun}
\label{Turing.RecordRun}
\leanok
\uses{Turing.Cfg, Turing.MultiTapeTM, Turing.MultiTapeTM.runFrom}
Let $R$ be a $K$-tape machine over input alphabet $\Gamma$ and state set $S$, let $x$ be
a finite input over $\Gamma$, let $a,b$ be configurations on $x$, let $y,n$ be
designated states, let $\beta$ assign an integer base position to each work tape, and
let $\rho,t$ be natural numbers. The run of $R$ from $a$ reaches $b$ in exactly $t$
steps; at every earlier time $u<t$ its state is neither $y$ nor $n$; and at every time
$u\le t$, each work-tape head is at distance at most $\rho$ from its assigned base
position.
\end{definition}

\begin{theorem}[Zero-Length Recorded Run]
\lean{Turing.recordRun_zero}
\label{Turing.recordRun_zero}
\leanok
\uses{Turing.Cfg, Turing.MultiTapeTM, Turing.RecordRun}
\difficulty{3}
Let $R$ be a multi-tape Turing machine, let $yes$ and $no$ be states, and choose a base
position for each work tape and a natural-number radius. If configuration $a$ has every
work-tape head within that radius of its base position, then the run from $a$ to itself
in zero steps avoids both distinguished states at every time strictly before its
endpoint and satisfies the head-position bound at every time through its endpoint.
\end{theorem}

\begin{theorem}[Concatenation of Recorded Runs]
\lean{Turing.recordRun_trans}
\label{Turing.recordRun_trans}
\leanok
\uses{Turing.Cfg, Turing.MultiTapeTM, Turing.MultiTapeTM.runFrom_add, Turing.RecordRun}
\difficulty{5}
Let $R$ be a multi-tape Turing machine with distinguished states $yes$ and $no$, and fix
a base head position for each work tape and a radius. If a run from configuration $a$
reaches $b$ in $t$ steps, avoids both distinguished states at every time before $t$, and
keeps every work-tape head within the specified radius of its base position through time
$t$, and a run from $b$ reaches $c$ in $v$ steps with the same avoidance and
head-position bounds, then the run from $a$ reaches $c$ in $t+v$ steps, avoids both
distinguished states before that time, and satisfies the same head-position bounds
throughout.
\end{theorem}

\begin{theorem}[State tag excludes a differently tagged state]
\lean{Turing.record_tag_ne}
\label{Turing.record_tag_ne}
\leanok
\uses{Turing.Cfg, Turing.Cfg.mapState}
\difficulty{4}
Let $K$ be a natural number, let $\Gamma$ and $I$ be types, and let $\mathrm{Ph}$ assign
a type to each element of $I$. Choose $i\in I$, an element $v$ of the dependent sum of
the types $\mathrm{Ph}(j)$ over $j\in I$, and a configuration $c$ on an input word $x$
whose states lie in $\mathrm{Ph}(i)$. If $i$ differs from the first component of $v$,
then after replacing each state $q$ by the pair $(i,q)$, the configuration's state is
not $v$.
\end{theorem}

\begin{theorem}[Tagged run satisfies recorded-run conditions]
\lean{Turing.recordRun_tagged}
\label{Turing.recordRun_tagged}
\leanok
\uses{Turing.Cfg, Turing.Cfg.mapState, Turing.MultiTapeTM, Turing.MultiTapeTM.runFrom, Turing.RecordRun, Turing.record_tag_ne}
\difficulty{5}
Let $R$ be a machine whose states are tagged phase states, and let $M$ be a machine with
phase states indexed by $i$. Suppose that through every time $u\le t$, the run of $R$
from $a$ agrees with the run of $M$ from $c$ after tagging its state by $i$. Suppose the
run of $R$ from $a$ reaches $b$ after $t+1$ steps, the phase index $i$ differs from the
phase indices of both distinguished states $\mathrm{yes}$ and $\mathrm{no}$, and the
work-head distances from the specified base positions are at most $r$ at every time
$u\le t$ in the run of $M$ and at time $t+1$ in configuration $b$. Then the run of $R$
from $a$ to $b$ in $t+1$ steps avoids both distinguished states at all earlier times and
has all work-head distances from the base positions at most $r$ at every time through
$t+1$.
\end{theorem}

\begin{theorem}[Width of the successor value]
\lean{Turing.recordStep_width_one}
\label{Turing.recordStep_width_one}
\leanok
\uses{Turing.posWidth}
\difficulty{4}
For every natural number $n$, the number of binary digits used to represent $n+1$ is one
exactly when $n=0$.
\end{theorem}

\begin{theorem}[Value of Complementary Binary Words]
\lean{Turing.recordStep_complement_value}
\label{Turing.recordStep_complement_value}
\leanok
\difficulty{4}
For every finite Boolean word $u$, interpret the word and its bitwise complement as
base-$2$ digits in their given order. The sum of their values is $2^{|u|}-1$.
\end{theorem}

\begin{theorem}[Binary word determined by its value]
\lean{Turing.recordStep_word_of_value}
\label{Turing.recordStep_word_of_value}
\leanok
\uses{Turing.binDigits, Turing.record_digits_value}
\difficulty{4}
Let $w$ be a finite Boolean word and $n\in\mathbb{N}$. If $n<2^{|w|}$ and interpreting
the entries of $w$ as binary digits, least significant digit first, gives the value $n$,
then $w$ is the length-$|w|$ list of binary digits of $n$, with each position recording
the corresponding binary place value.
\end{theorem}

\begin{theorem}[Successor and Predecessor of Fixed-Width Digits]
\lean{Turing.recordStep_digits_succ}
\label{Turing.recordStep_digits_succ}
\leanok
\uses{Turing.binDigits, Turing.counterWord, Turing.counterWord_length, Turing.counterWord_value, Turing.decFixed, Turing.incFixed, Turing.recordStep_word_of_value, Turing.record_digits_inj, Turing.record_digits_value}
\difficulty{6}
Let $w,n$ be natural numbers with $n+1<2^w$. The fixed-width little-endian binary
increment of the $w$-digit representation of $n$ is the representation of $n+1$, and the
fixed-width decrement of that representation is the representation of $n$.
\end{theorem}

\begin{theorem}[Saved counter digits]
\lean{Turing.recordStep_saved_digits}
\label{Turing.recordStep_saved_digits}
\leanok
\uses{Turing.binDigits, Turing.counterWord, Turing.counterWord_length, Turing.counterWord_value, Turing.recordStep_complement_value, Turing.recordStep_word_of_value}
\difficulty{3}
Let $w$ be a Boolean word and $q$ a natural number. If the binary value of $w$ is
$2^{|w|}-1$ and $q+1<2^{|w|}$, then complementing each bit of the word obtained by
applying the fixed-width decrement operation $q+1$ times to $w$ gives the width-$|w|$
binary digit list of $q+1$.
\end{theorem}

\begin{theorem}[Preservation of the tape window]
\lean{Turing.recordStep_window}
\label{Turing.recordStep_window}
\leanok
\uses{Turing.Action, Turing.Action.apply, Turing.Cfg, Turing.Cfg.InWindow}
\difficulty{5}
Let $c$ be a configuration of a machine with $k$ Boolean work tapes whose heads lie in
the integer interval $[-s,s]$ and whose tape contents are blank outside that interval,
where $s\in\mathbb{N}$. For any action $a$, the configuration obtained by applying $a$
still has all heads in that interval and all tape contents blank outside it if and only
if, on every tape, the head is not at $-s$ while moving left and is not at $s$ while
moving right.
\end{theorem}

\begin{theorem}[Restriction to embedded Boolean tapes]
\lean{Turing.recordStep_restrict}
\label{Turing.recordStep_restrict}
\leanok
\uses{Turing.Cfg, Turing.Cfg.liftSymbols, Turing.MultiTapeTM, Turing.MultiTapeTM.liftSymbols, Turing.decodeBool, Turing.decodeBool_apply, Turing.embedEmitCfg, Turing.embedEmitCfg_selected_pos, Turing.embedEmitCfg_selected_tape, Turing.embedEmitTM_frame}
\difficulty{5}
Let $K,m$ be natural numbers, let $\Gamma,S$ be types with decidable equality, let $R$
be an $m$-work-tape Boolean Turing machine with states in $S$, let
$e:\{\mathrm{false},\mathrm{true}\}\hookrightarrow\Gamma$ be an embedding, and let
$\iota:\{0,\ldots,m-1\}\hookrightarrow\{0,\ldots,K-1\}$ be an embedding. For a Boolean
input word $x$, consider a $K$-tape configuration $d$ over $\Gamma$ on the input
obtained by applying $e$ to every bit of $x$. Assume every symbol stored anywhere on
each tape indexed by $\iota$ lies in the image of $e$. Then $d$ is exactly the
configuration formed by embedding the decoded Boolean configuration into the selected
tapes via $\iota$, leaving all other tapes and head positions as in $d$, and prefixing
the decoded configuration's output by $d$'s output.
\end{theorem}

\begin{theorem}[Transport under symbol and tape embeddings]
\lean{Turing.recordStep_transport}
\label{Turing.recordStep_transport}
\leanok
\uses{Turing.Cfg, Turing.Cfg.liftSymbols, Turing.MultiTapeTM, Turing.MultiTapeTM.liftSymbols, Turing.MultiTapeTM.liftSymbols_runFrom, Turing.MultiTapeTM.runFrom, Turing.decodeBool, Turing.embedEmitCfg, Turing.embedEmitTM, Turing.embedEmitTM_runFrom, Turing.recordStep_restrict}
\difficulty{1}
Let $R$ be a deterministic machine with $m$ Boolean work tapes and state set $S$, let
$e$ embed the Boolean alphabet into an alphabet $\Gamma$, and let $\iota$ embed the $m$
tape indices into $K$ tape indices. Suppose $d$ is a configuration on input obtained by
applying $e$ to each symbol of a Boolean word $x$, and every nonblank symbol on each
embedded tape $\iota(i)$ belongs to the image of $e$. Decode those embedded tapes using
$e$, run $R$ for $t\in\mathbb{N}$ steps, and then lift its symbols by $e$. The resulting
configuration, with the original embedded tapes and head positions retained on unused
tape indices and the original output prefixed to the lifted output, is exactly the
configuration reached after $t$ steps by the machine obtained from $R$ by lifting its
symbols and embedding its tapes.
\end{theorem}

\begin{theorem}[Preservation of the work-tape frame]
\lean{Turing.recordStep_native_frame}
\label{Turing.recordStep_native_frame}
\leanok
\uses{Turing.Cfg, Turing.MultiTapeTM, Turing.MultiTapeTM.runFrom, Turing.MultiTapeTM.runFrom_succ_eq_step', Turing.MultiTapeTM.step_of_halt, Turing.inputMoveAction, Turing.step_of_inputMoveAction}
\difficulty{4}
Let $R$ be a $K$-tape machine over an alphabet $\Gamma$ with state set $S$, and suppose
its transition on every state, input symbol, and tuple of scanned work-tape symbols is
the action that moves only the input head and changes the state according to functions
$\mathrm{move}:S\to\mathrm{SignType}$ and $\mathrm{next}:S\to\mathrm{Option}(\Gamma)\to
S$. Then, from any configuration $d$ and after any number $t$ of steps, the work-tape
contents, work-head positions, and output are unchanged.
\end{theorem}

\begin{definition}[Record step probe states]
\lean{Turing.RecordStepProbe}
\label{Turing.RecordStepProbe}
\leanok
The record step probe has three states: a start state, a peek state, and a done state
carrying a Boolean value.
\end{definition}

\begin{definition}[Record step probe machine]
\lean{Turing.recordStepProbe}
\label{Turing.recordStepProbe}
\leanok
\uses{Turing.MultiTapeTM, Turing.RecordStepProbe, Turing.inputMoveAction}
For any symbol alphabet, a four-work-tape machine has the start state, with states given
by the record step probe states. In the start state, if the input head scans a symbol,
it stays in place and moves to the done state with value false; if the input is blank,
it moves right and enters the peek state. In the peek state it moves left and enters the
done state with value true exactly when the input is blank, while each done state stays
in place and remains unchanged.
\end{definition}

\begin{definition}[Record step save machine]
\lean{Turing.recordStepSave}
\label{Turing.recordStepSave}
\leanok
\uses{Turing.MultiTapeTM, Turing.inputMoveAction}
A machine with no work tapes, Boolean input alphabet, and Boolean states starts in state
false. It stays in place and enters state true if its current state is true or the input
is blank; otherwise, it moves left and enters state false.
\end{definition}

\begin{definition}[Record step read machine]
\lean{Turing.recordStepRead}
\label{Turing.recordStepRead}
\leanok
\uses{Turing.MultiTapeTM, Turing.RecSym, Turing.RecSym.digitEmb, Turing.decodeBool, Turing.inputMoveAction}
Let $Q$ and $T$ be types with decidable equality. A machine with four work tapes,
alphabet consisting of record symbols over $Q$ and $T$, and state set
$\text{Option}(\text{Option}(\text{Bool}))$ starts in the state indicating halt. At every
step it stays in place and enters the state obtained by decoding the scanned input
symbol when it is a binary digit; a blank or any non-digit input gives the inner value
none, and the resulting value is wrapped as a nonhalting state.
\end{definition}

\begin{theorem}[One-Step Read of a Boolean Digit]
\lean{Turing.recordStep_read_run}
\label{Turing.recordStep_read_run}
\leanok
\uses{Turing.Cfg, Turing.Cfg.inputSymbol, Turing.MultiTapeTM.runFrom, Turing.MultiTapeTM.step, Turing.RecSym, Turing.RecSym.digitEmb, Turing.decodeBool, Turing.moveInputPos_zero, Turing.recordStepRead, Turing.step_of_inputMoveAction}
\difficulty{3}
Let $Q$ and $T$ be types with decidable equality, let $x$ be a finite word over the
record alphabet determined by $Q$ and $T$, and let $d$ be a four-work-tape configuration
of the digit-reading machine on input $x$. If the state of $d$ is the initial state,
then after one step the configuration is exactly $d$ with its state replaced by the
state that records the optional Boolean decoded from the symbol under the input head (a
digit symbol decodes to its Boolean value, and every other symbol decodes to no value).
Moreover, at every time $u<1$, the machine's state is not any state of the form that
records a Boolean value.
\end{theorem}

\begin{theorem}[Input-head orbit of the saving machine]
\lean{Turing.recordStep_save_orbit}
\label{Turing.recordStep_save_orbit}
\leanok
\uses{Turing.Cfg, Turing.Cfg.inputSymbol, Turing.FinTM.moveInputPos_neg_val, Turing.MultiTapeTM.step, Turing.counterLoopOrbit, Turing.inputSymbolInner, Turing.moveInputPos, Turing.moveInputPos_zero, Turing.recordStepSave, Turing.roundOrbit, Turing.step_of_inputMoveAction}
\difficulty{6}
Let $x$ be a finite Boolean word, and let $c$ be a configuration of the saving machine
on $x$ with no work tapes. Suppose the machine is in state false and the input-head
position is at most the length of $x$. For every natural number $r$ no greater than that
position, after $r$ rounds—each consisting of one machine step—the state remains false,
the input-head position is its initial position minus $r$, and the output is unchanged.
After one more round than the initial input-head position, the state is true, the
input-head position is zero, and the output is unchanged.
\end{theorem}

\begin{theorem}[Characterization of Blank Input Positions]
\lean{Turing.recordStep_input_blank}
\label{Turing.recordStep_input_blank}
\leanok
\uses{Turing.Cfg, Turing.Cfg.inputSymbol, Turing.inputSymbolInner}
\difficulty{4}
Let $x$ be a finite word and let $d$ be a configuration with input-head position among
the $|x|+2$ available positions, including the two boundary blanks. The input symbol
read by $d$ is absent if and only if the head is at position $0$ or at position $|x|+1$.
\end{theorem}

\begin{theorem}[Input-boundary probe run]
\lean{Turing.recordStep_probe_run}
\label{Turing.recordStep_probe_run}
\leanok
\uses{Turing.Cfg, Turing.Cfg.inputSymbol, Turing.FinTM.moveInputPos_neg_val, Turing.MultiTapeTM.runFrom, Turing.MultiTapeTM.step, Turing.RecordStepProbe, Turing.inputSymbolInner, Turing.moveInputPos, Turing.moveInputPos_rightBoundary, Turing.moveInputPos_zero, Turing.recordStepProbe, Turing.recordStep_input_blank, Turing.recordStep_native_frame, Turing.step_of_inputMoveAction}
\difficulty{5}
Let $\Gamma$ be a type, let $x$ be a word over $\Gamma$ with positive length, and let
$d$ be a four-work-tape machine configuration on input $x$ whose state is the probe's
start state. Then there exist a time $t\le 2$, a Boolean value $b$, and an input-head
position $p$ such that after $t$ steps the probe is in state done with flag $b$ and
input-head position $p$. Before time $t$ it is in no done state; moreover, $b$ is true
exactly when the initial input-head position is the right boundary $|x|+1$, $p$ equals
$|x|$ when $b$ is true and otherwise equals the initial input-head position, and all
work-tape head positions remain unchanged through time $t$.
\end{theorem}

\begin{theorem}[Bounded counter-saving run]
\lean{Turing.recordStep_save_run}
\label{Turing.recordStep_save_run}
\leanok
\uses{Turing.Cfg, Turing.Cfg.ext_zero_tapes, Turing.FinTM.bufferTape, Turing.MultiTapeTM.runFrom, Turing.counterLoopCfg, Turing.counterLoopOrbit, Turing.counterLoopTM, Turing.counterLoopTM_run_escape, Turing.counterOverhead, Turing.counterOverhead_le_of_le, Turing.counterWord, Turing.counterWord_length, Turing.counterWord_value, Turing.recordStepSave, Turing.recordStep_complement_value, Turing.recordStep_save_orbit}
\difficulty{4}
Let $x$ and $w$ be finite Boolean words, and let $c$ be a configuration on input $x$
with no work tapes. Suppose its state is false and its input-head position $p$ satisfies
$p\le |x|$. Assume that the binary value of $w$ is $2^{|w|}-1$ and that $p$ is strictly
less than this value. Let $ct:\mathbb Z\to\mathrm{Option}(\mathrm{Bool})$ be the
contents of an additional work tape, with head position $cp\in\mathbb Z$, and suppose
that for every integer $z$ with $-1\le z\le |w|$, the cell at position $cp+z$ contains
the symbol of $w$ at offset $z$ (or is blank when that offset is outside the word).
Consider the counter-loop machine whose body, on a false state reading an input symbol,
moves the input head one position left and remains false, and which otherwise enters the
true exit state; its counter is stored on the additional tape and decremented between
body iterations. There is a time $t\le 5(p+1)+2|w|$ at which this machine is in its
escape state, its embedded configuration has true state and input-head position zero,
and its additional tape contains the word obtained by applying the fixed-width decrement
operation to $w$ exactly $p+1$ times. Before time $t$ it is in neither the done nor
escape state, and at every time through $t$ the additional-tape head lies in
$[cp-1,cp+|w|]$. The binary value of the bitwise complement of the resulting word is
$p+1$.
\end{theorem}

\begin{definition}[Record step data]
\lean{Turing.RecordStepData}
\label{Turing.RecordStepData}
\leanok
Record-step data for $k$ work tapes, a state type $Q$, and a tag type $T$ consists of an
optional state, a Boolean width flag, a Boolean position flag, an optional Boolean
symbol for each tape, left-edge and right-edge Boolean flags for each tape, an optional
tag, a Boolean native-input direction, and an optional Boolean input symbol. By default,
the state, tape symbols, tag, and input symbol are absent; the width and
native-direction flags are false, the position flag is true, and all edge flags are
false.
\end{definition}

\begin{definition}[Record step stages]
\lean{Turing.RecordStepStage}
\label{Turing.RecordStepStage}
\leanok
For $k$ work tapes, the record-step stages include state handling, position handling,
copying, tape reading and skipping for each tape, tag reading, reading back and writing
back indexed fields, halt and empty-cell clearing, fuel and probe operations, saving and
seeking, reading and complementing, counter updates, restoration, state and tape
writing, and verdict stages carrying a Boolean value. The indexed read-back and
write-back stages range over $k+2$ positions.
\end{definition}

\begin{definition}[Record step index]
\lean{Turing.RecordStepIndex}
\label{Turing.RecordStepIndex}
\leanok
\uses{Turing.RecordStepData, Turing.RecordStepStage}
For a natural number $k$ and types $Q$ and $T$, the record-step index is the Cartesian
product of the record-step stages for $k$ tapes and the corresponding record-step data.
\end{definition}

\begin{definition}[Record Step Tape Fold]
\lean{Turing.recordStepTapeFold}
\label{Turing.recordStepTapeFold}
\leanok
\uses{Turing.RecSym}
For types $Q$ and $T$, define a function from $\text{Bool} \times
\text{Option}(\text{Bool}) \times \text{Bool} \times \text{Bool}$ and record symbols to
the same product as follows. On a cell symbol with content $u$ and head mark $m$, an
input accumulator $(a,v,c,d)$ is sent to $(\mathrm{true},\, v \text{ if }
m=\mathrm{false} \text{ and } u \text{ if } m=\mathrm{true},\, c \lor (\neg a \land
m),\, m)$; every other record symbol leaves the accumulator unchanged.
\end{definition}

\begin{theorem}[Tape fold on a finite prefix]
\lean{Turing.recordStep_tape_prefix}
\label{Turing.recordStep_tape_prefix}
\leanok
\uses{Turing.recordStepTapeFold}
\difficulty{5}
Let $Q,T$ be types, let $f:\mathbb{N}\to\operatorname{Option}(\mathrm{Bool})$, and let
$p,r$ be natural numbers. Folding the first $r$ record symbols, with symbol $j$
containing $f(j)$ and a mark exactly when $p=j$, from the initial summary
$(\mathrm{false},\mathrm{none},\mathrm{false},\mathrm{false})$ yields $( [0<r],\,
f(p)\text{ if }p<r\text{ and none otherwise},\, [0<r\text{ and }p=0],\, [0<r\text{ and
}p+1=r])$, where each bracket denotes the Boolean truth value of the stated condition.
\end{theorem}

\begin{theorem}[Fold over a tape window]
\lean{Turing.recordStep_tape_fold}
\label{Turing.recordStep_tape_fold}
\leanok
\uses{Turing.Cfg, Turing.Cfg.InWindow, Turing.RecSym, Turing.recordStepTapeFold, Turing.recordStep_tape_prefix, Turing.scanFold, Turing.scanKeep, Turing.scanResult, Turing.scanResult_fold}
\difficulty{5}
Let $Q,T$ be types, let $s$ be a natural number, and let $c$ be a configuration of a
$k$-tape Boolean machine on input $x$. Suppose every work-tape head of $c$ lies within
distance $s$ of the origin and every tape is blank at integer positions whose absolute
value exceeds $s$. For any work-tape index $i$, form the record-symbol list indexed by
$j=0,\ldots,2s$, whose $j$th entry stores the content of tape $i$ at position $j-s$ and
marks whether its head is there. Scanning this list with the tape fold in the forward
direction and then the identity backward leaves the list unchanged and returns the
summary $(\mathrm{true},c_i(h_i),[h_i=-s],[h_i=s])$, where $h_i$ is the head position,
$c_i(h_i)$ is the optional symbol under that head, and each bracket is the Boolean truth
value of the equality.
\end{theorem}

\begin{definition}[Digit position fold]
\lean{Turing.recordStepPosFold}
\label{Turing.recordStepPosFold}
\leanok
\uses{Turing.RecSym}
For a symbol that is a binary digit, update a pair consisting of a position in
$\{0,1,2\}$ and a Boolean flag by replacing the position with the smaller of $2$ and one
plus its current value, and replacing the flag with its conjunction with the negation of
the digit. For every other symbol, leave the pair unchanged.
\end{definition}

\begin{theorem}[Position fold over binary digits]
\lean{Turing.recordStep_position_fold}
\label{Turing.recordStep_position_fold}
\leanok
\uses{Turing.recordStepPosFold, Turing.scanFold, Turing.scanKeep, Turing.scanResult, Turing.scanResult_fold}
\difficulty{5}
For every Boolean word $w$, scanning its symbols as recording-alphabet digit symbols
with the position fold, starting from count zero and an initially true flag, returns the
original word unchanged together with count $\min(2,|w|)$ and the assertion that the
natural number encoded by $w$ in base two is zero.
\end{theorem}

\begin{definition}[Recorded machine action]
\lean{Turing.recordStepAction}
\label{Turing.recordStepAction}
\leanok
\uses{Turing.Action, Turing.MultiTapeTM, Turing.RecordStepData}
Given a deterministic $k$-tape Turing machine over the binary alphabet and record data,
return the machine's transition action when the recorded state is present, using the
recorded input and work-tape symbols. If the recorded state is absent, return an action
that stays on the input and work tapes, writes nothing, emits nothing, and halts.
\end{definition}

\begin{definition}[Recorded work-tape boundary exit]
\lean{Turing.recordStepOutside}
\label{Turing.recordStepOutside}
\leanok
\uses{Turing.MultiTapeTM, Turing.RecordStepData, Turing.recordStepAction}
Given a deterministic $k$-tape Turing machine over the binary alphabet and record data,
the returned Boolean is true exactly when some work tape is marked as being at its left
edge and its action moves left, or is marked as being at its right edge and its action
moves right; otherwise it is false.
\end{definition}

\begin{definition}[Bank selection for seeking]
\lean{Turing.recordStepSeekBank}
\label{Turing.recordStepSeekBank}
\leanok
For each Boolean choice of whether to restore, define an injective map from a
two-element set to a four-element set: the first element maps to the fourth element when
restoring and to the second element otherwise, while the second element maps to the
third element.
\end{definition}

\begin{definition}[Bank selection for digit records]
\lean{Turing.recordStepDigitBank}
\label{Turing.recordStepDigitBank}
\leanok
For each $i\in\{0,1,2\}$, define an injective map from a one-element set to a
four-element set that sends its unique element to the element indexed by $i+1$.
\end{definition}

\begin{definition}[Extracting a cell mark]
\lean{Turing.recordStepGetMark}
\label{Turing.recordStepGetMark}
\leanok
\uses{Turing.RecSym}
For a recursive record symbol, return its mark if it represents a work cell, and return
false for every other kind of symbol.
\end{definition}

\begin{definition}[Updating a cell mark]
\lean{Turing.recordStepSetMark}
\label{Turing.recordStepSetMark}
\leanok
\uses{Turing.RecSym}
Given a recursive record symbol and a Boolean value, replace the mark with that value if
the symbol represents a work cell, preserving its content. Leave every other kind of
symbol unchanged.
\end{definition}

\begin{definition}[Writing to a marked cell]
\lean{Turing.recordStepWriteCell}
\label{Turing.recordStepWriteCell}
\leanok
\uses{Turing.RecSym}
Given an optional write instruction and a recursive record symbol, update the content
only when the symbol represents a work cell whose mark is true. In that case, use the
instructed content when present and otherwise preserve the existing content; preserve
the mark, and leave every other kind of symbol unchanged.
\end{definition}

\begin{theorem}[Right boundary detection]
\lean{Turing.recordStep_right_clamp}
\label{Turing.recordStep_right_clamp}
\leanok
\uses{Turing.Cfg, Turing.Cfg.inputSymbol, Turing.recordStep_input_blank}
\difficulty{4}
For any natural number $k$, type $Q$, Boolean input word $x$, and configuration $c$ of a
$k$-tape machine on $x$, the conjunction that the input head reads no symbol and is not
at position zero is equivalent to the input head being at position $|x|+1$.
\end{theorem}

\begin{theorem}[Binary digits track input-head movement]
\lean{Turing.recordStep_position_digits}
\label{Turing.recordStep_position_digits}
\leanok
\uses{Turing.Cfg, Turing.Cfg.inputSymbol, Turing.FinTM.moveInputPos_neg_val, Turing.binDigits, Turing.decFixed, Turing.incFixed, Turing.moveInputPos, Turing.moveInputPos_pos_of_ne_right, Turing.moveInputPos_rightBoundary, Turing.posWidth, Turing.recordStep_digits_succ, Turing.recordStep_right_clamp, Turing.record_pos_lt}
\difficulty{5}
Let $c$ be a configuration of a machine on a Boolean input $x$, and let $m$ be a left,
stationary, or right move. Encode the input-head position using fixed-width,
little-endian binary digits of width equal to the binary width of $|x|+1$. If $m$ is
stationary, the digits are unchanged; if $m$ is left, they are unchanged at the left
boundary and otherwise undergo fixed-width decrement; if $m$ is right, they are
unchanged when the scanned input symbol is blank and the head is not at the left
boundary, and otherwise undergo fixed-width increment. In every case, the resulting word
exists and is exactly the fixed-width binary encoding of the input-head position after
the clamped move.
\end{theorem}

\begin{theorem}[Recorded Work-Cell Update]
\lean{Turing.recordStep_cell_write}
\label{Turing.recordStep_cell_write}
\leanok
\uses{Turing.Action, Turing.Action.apply, Turing.Cfg, Turing.recordStepWriteCell}
\difficulty{4}
Let $a$ be an action for a Boolean multi-tape machine, let $c$ be a configuration, and
let $i$ be a work-tape index and $z$ an integer cell position. The recorded cell whose
content is that of tape $i$ at $z$ and whose mark indicates whether the head is at $z$
is transformed by the action's write instruction for tape $i$ into a cell whose content
is exactly the content of tape $i$ at $z$ after applying $a$ to $c$, while the mark
remains unchanged.
\end{theorem}

\begin{theorem}[Recorded Cell Mark Shift]
\lean{Turing.recordStep_cell_shift}
\label{Turing.recordStep_cell_shift}
\leanok
\uses{Turing.RecSym, Turing.recordStepGetMark, Turing.recordStepSetMark, Turing.scanKeep, Turing.scanResult, Turing.scanResult_fold, Turing.scanResult_shiftLeft, Turing.scanResult_shiftRight, Turing.scanShift}
\difficulty{5}
Let $s$ be a natural number, let $f$ assign an optional Boolean value to each integer,
and let $h$ be an integer with $|h|\le s$. For each movement $m\in\{-1,0,1\}$, form the
ordered list of $2s+1$ cells indexed by $j-s$ for $j=0,\ldots,2s$, with contents
$f(j-s)$ and a mark precisely at $h$. The specified scan leaves every cell’s content
unchanged and marks precisely the cell indexed by $h+m$.
\end{theorem}

\begin{theorem}[Recorded tape word after an action]
\lean{Turing.recordStep_tape_word}
\label{Turing.recordStep_tape_word}
\leanok
\uses{Turing.Action, Turing.Action.apply, Turing.Cfg, Turing.Cfg.InWindow, Turing.RecSym, Turing.recordStepGetMark, Turing.recordStepSetMark, Turing.recordStepWriteCell, Turing.recordStep_cell_shift, Turing.recordStep_cell_write, Turing.scanKeep, Turing.scanMap, Turing.scanResult, Turing.scanResult_map, Turing.scanShift}
\difficulty{3}
Let $c$ be a configuration of a machine with $k$ Boolean work tapes, let its work-tape
heads lie within distance $s\in\mathbb{N}$ of the origin and its tape contents be blank
outside that window, and let $a$ be an action and $i<k$. Form the word of $2s+1$
recorded cells indexed from $-s$ to $s$, storing each cell's content and a mark for
whether the head of tape $i$ is there. First apply the action's prescribed write to the
marked cell; then shift the mark one place left, leave it fixed, or shift it one place
right according as the action moves that tape's head left, keeps it still, or moves it
right. The resulting word is exactly the $2s+1$ recorded cells of tape $i$ in the
configuration obtained by applying $a$ to $c$, with marks at the new head position.
\end{theorem}

\begin{definition}[Phase type for recording a step]
\lean{Turing.recordStepPhase}
\label{Turing.recordStepPhase}
\leanok
\uses{Turing.CounterLoopState, Turing.FlagPhase, Turing.RecordStepIndex, Turing.RecordStepProbe, Turing.ScanPhase, Turing.SeekPhase, Turing.SweepPhase}
For a natural number $k$, types $Q,T,S$, and a record-step index consisting of a stage
and its associated data, assign a phase type according to the stage. The state,
position, tape-reading, tag-reading, writing, and shifting stages use scan phases
carrying, respectively, an optional state, a pair in $\text{Fin}(3)\times\mathrm{Bool}$,
a tuple in
$\text{Bool}\times\mathrm{Option}(\text{Bool})\times\mathrm{Bool}\times\mathrm{Bool}$, an
optional tag, the unit type, or a Boolean; the probe and save stages use the probe and
counter-loop state types, the seek and restore stages use $S$, the read stage uses a
nested optional Boolean, the counter-update stages use flag phases, and the copying and
clearing stages use sweep phases. Every remaining stage uses a seek phase.
\end{definition}

\begin{definition}[Record-step machine construction]
\lean{Turing.recordStepRows}
\label{Turing.recordStepRows}
\leanok
\uses{Turing.MultiTapeTM, Turing.MultiTapeTM.liftSymbols, Turing.RecSym, Turing.RecSym.digitEmb, Turing.RecordStepIndex, Turing.clearTM, Turing.copyTM, Turing.counterLoopTM, Turing.decrementTM, Turing.embedEmitTM, Turing.incrementTM, Turing.recordStepAction, Turing.recordStepDigitBank, Turing.recordStepGetMark, Turing.recordStepPhase, Turing.recordStepPosFold, Turing.recordStepProbe, Turing.recordStepRead, Turing.recordStepSave, Turing.recordStepSeekBank, Turing.recordStepSetMark, Turing.recordStepTapeFold, Turing.recordStepWriteCell, Turing.scanFold, Turing.scanKeep, Turing.scanMap, Turing.scanShift, Turing.scanTM, Turing.skipBackTM, Turing.skipTM}
Let $k$ be a natural number, let $Q$, $T$, and $S$ be types with decidable equality on
$Q$ and $T$, let $M$ be a deterministic Turing machine with $k$ work tapes, Boolean
symbols, and state set $Q$, and let $u:T\to\mathrm{Option}(\mathrm{Bool})\to T$ be a
tag-update function. Let $N$ be a deterministic Turing machine with two work tapes,
Boolean symbols, and state set $S$, and let $i$ be a record-step index for $k$, $Q$, and
$T$. The construction assigns to these data a deterministic Turing machine with four
work tapes, record symbols formed from $Q$ and $T$, and control phases determined by
$i$; its transition behavior is defined by cases on the stage and stored record data.
\end{definition}

\begin{definition}[Initial Control for a Recorded Step]
\lean{Turing.recordStepApply}
\label{Turing.recordStepApply}
\leanok
\uses{Turing.MultiTapeTM, Turing.RecordStepData, Turing.RecordStepIndex, Turing.recordStepOutside, Turing.recordStepPhase}
Let $k$ be a natural number, let $Q$, $T$, and $S$ be types, let $M$ be a deterministic
Turing machine with $k$ work tapes and Boolean input symbols, and let $d$ be record-step
data for $Q$ and $T$. The construction returns a record-step index together with its
phase: if the record lies outside the permitted range, the stage is clear-reject with
sweep phase; otherwise, the stage is write-state with forward scan phase carrying the
unit value.
\end{definition}

\begin{definition}[First tape stage selection]
\lean{Turing.recordStepFirstTape}
\label{Turing.recordStepFirstTape}
\leanok
\uses{Turing.RecordStepData, Turing.RecordStepIndex, Turing.recordStepPhase}
For a natural number $k$, types $Q,T,S$, record-step data, and a Boolean indicating
whether the stage is writing, select the first work-tape stage when $k>0$: use the write
stage if writing, and otherwise the tape-reading stage with its initial scan data. When
$k=0$, use the tag-writing stage if writing, and otherwise the tag-reading stage with no
tag initially recorded.
\end{definition}

\begin{definition}[Successor control for recording stages]
\lean{Turing.recordStepNext}
\label{Turing.recordStepNext}
\leanok
\uses{Turing.MultiTapeTM, Turing.RecordStepIndex, Turing.recordStepAction, Turing.recordStepApply, Turing.recordStepFirstTape, Turing.recordStepPhase}
Given a natural number $k$, types $Q,T,S$, a Boolean-symbol multi-tape Turing machine
$M$ with $k$ work tapes, start and finish states in $S$, a record-step index, and
decidable equality on $S$, define a partial successor function from the phase associated
with that index to an optional next index-phase pair. The successor is determined by the
stage and phase cases: completed scans and sweeps advance through the recording
procedure, update the stored data when required, and may select a verdict; the seek and
restore stages continue only at the specified finish state, and phases not matching a
case have no successor.
\end{definition}

\begin{definition}[Control type for recording a step]
\lean{Turing.RecordStepControl}
\label{Turing.RecordStepControl}
\leanok
\uses{Turing.RecordStepIndex, Turing.recordStepPhase}
For a natural number $k$ and types $Q,T,S$, the record-step control type consists of a
record-step index together with a phase belonging to the phase type assigned to that
index.
\end{definition}

\begin{definition}[Entry control for recording a step]
\lean{Turing.recordStepEntry}
\label{Turing.recordStepEntry}
\leanok
\uses{Turing.RecordStepControl}
For a natural number $k$ and types $Q,T,S$, the entry control is the state-reading stage
with default record-step data and forward scan phase carrying no state.
\end{definition}

\begin{definition}[Verdict control for recording a step]
\lean{Turing.recordStepVerdict}
\label{Turing.recordStepVerdict}
\leanok
\uses{Turing.RecordStepControl}
For a natural number $k$, types $Q,T,S$, and a Boolean verdict, the verdict control
consists of the verdict stage carrying that Boolean, default record-step data, and the
run phase.
\end{definition}

\begin{definition}[Record step machine]
\lean{Turing.recordStepMachine}
\label{Turing.recordStepMachine}
\leanok
\uses{Turing.MultiTapeTM, Turing.RecSym, Turing.RecordStepControl, Turing.phaseHostTM, Turing.recordStepEntry, Turing.recordStepNext, Turing.recordStepRows}
Given a natural number $k$, a Boolean Turing machine $M$ with $k$ work tapes and state
set $Q$, a function $u:T\times\mathrm{Option}(\mathrm{Bool})\to T$, a Boolean Turing
machine $S$ with two work tapes and state set $S$, and designated states $s_0,s_f\in S$,
the definition constructs a Turing machine with four work tapes, record symbols built
from $Q$ and $T$, and control states encoding its processing stage and associated data
and phase. The resulting machine is specified by the record-step transition rows, entry
state, and next-state function determined by $M,u,S,s_0,s_f$; assume that $Q,T,S$ each
have decidable equality.
\end{definition}

\begin{theorem}[State scan in five steps]
\lean{Turing.recordStep_state_run}
\label{Turing.recordStep_state_run}
\leanok
\uses{Turing.Cfg, Turing.Cfg.mapState, Turing.FinTM.bufferTape, Turing.MultiTapeTM, Turing.MultiTapeTM.runFrom, Turing.RecSym, Turing.RecordRun, Turing.RecordStepControl, Turing.RecordStepIndex, Turing.ScanPhase, Turing.phaseHostTM_run, Turing.recordRun_tagged, Turing.recordStepEntry, Turing.recordStepMachine, Turing.recordStepNext, Turing.recordStepRows, Turing.recordStepVerdict, Turing.scanFold, Turing.scanKeep, Turing.scanResult, Turing.scanResult_fold, Turing.scanTM, Turing.scanTM_run_ofCfg}
\difficulty{5}
Let $Q,T,S$ be types with decidable equality. Let $M$ be a Boolean multi-tape machine
with $k$ work tapes, let
$\operatorname{upd}:T\times\operatorname{Option}(\mathrm{Bool})\to T$, and let
$\mathrm{seek}$ be a two-work-tape Boolean machine with designated states
$\mathrm{start},\mathrm{finish}$. For any input list $x$ of record symbols, four-tape
configuration $d$, and optional state $q:Q$, suppose $d$ is in the record-step entry
control state and tape 0, at offsets $-1,0,1$ from its head, contains the one-symbol
word consisting of the state symbol encoding $q$. After five steps of the record-step
machine, the configuration is $d$ with its control state changed to the running phase
that skips the state symbol and carries record data whose stored state is $q$; neither
verdict state has been reached before those five steps, and every work-tape head remains
within distance one of its initial position through that time.
\end{theorem}

\begin{theorem}[One phase transition in the hosted run]
\lean{Turing.recordStep_phase_run}
\label{Turing.recordStep_phase_run}
\leanok
\uses{Turing.Cfg, Turing.Cfg.mapState, Turing.MultiTapeTM, Turing.MultiTapeTM.runFrom, Turing.RecSym, Turing.RecordRun, Turing.RecordStepControl, Turing.RecordStepIndex, Turing.phaseHostTM_run, Turing.recordRun_tagged, Turing.recordStepEntry, Turing.recordStepMachine, Turing.recordStepNext, Turing.recordStepPhase, Turing.recordStepRows, Turing.recordStepVerdict}
\difficulty{3}
Let $M$ be a $k$-tape Boolean machine with state set $Q$, let
$\mathrm{upd}:T\to\mathrm{Option}(\mathrm{Bool})\to T$, and let $\mathrm{seek}$ be a
two-tape Boolean machine with designated states $\mathrm{start},\mathrm{finish}\in S$.
Choose a record-stage index $i$ whose stage is not a verdict stage, and configurations
$c,c'$ of the corresponding phase machine on an input $x$. Suppose that running its row
machine from $c$ for $t$ steps gives $c'$, that $c'$ has phase state $s'$, and that the
phase-transition function sends $s'$ to an outcome $o$. Suppose also that before time
$t$ every phase state reached has no phase transition, and that through time $t$ every
work-head position stays within radius $r$ of a specified base position on each of the
four tapes. Then the hosted record-step machine runs from the configuration obtained by
adjoining phase index $i$ to $c$ to the configuration obtained by replacing the state of
$c'$ with $o$, in $t+1$ steps, without visiting either verdict state before the final
time; throughout these steps all four work-head displacements from the specified base
positions are at most $r$.
\end{theorem}

\begin{theorem}[Preservation during a two-pass scan]
\lean{Turing.recordStep_scan_preserves}
\label{Turing.recordStep_scan_preserves}
\leanok
\uses{Turing.Cfg, Turing.FinTM.bufferTape, Turing.MultiTapeTM.runFrom, Turing.ScanPhase, Turing.scanResult, Turing.scanTM, Turing.scanTM_run_ofCfg}
\difficulty{4}
Let $A$ and $\Gamma$ be types, let $x,w$ be lists over $\Gamma$, let $i$ be one of four
work-tape indices, and let $\delta_f,\delta_b:A\times\Gamma\to A\times\Gamma$ be
transition maps. Suppose a configuration $c$ of the four-tape scan machine on input $x$
is in forward phase with state $a$, and tape $i$, at every offset $z$ from its head with
$-1\le z\le |w|$, contains the corresponding symbol of $w$ (and is blank where that word
has no symbol). If the forward-and-backward scan of $w$ from state $a$ returns the pair
$(b,w)$, then after $2|w|+2$ steps the machine is in configuration $c$ with state
changed to the completed phase carrying $b$. It has not entered any completed phase
earlier, and through that time tape $i$'s head remains between its initial position
minus one and its initial position plus $|w|$, while every other work-tape head stays
fixed.
\end{theorem}

\begin{definition}[Record Step Snapshot]
\lean{Turing.recordStepSnapshot}
\label{Turing.recordStepSnapshot}
\leanok
\uses{Turing.Cfg, Turing.FinTM.bufferTape}
Let $d$ be a configuration on input $x$ with four work tapes, state set $P$, and symbol
alphabet $\Gamma$. For a new state $q\in A$, integer $p$, and lists $b,e$ over $\Gamma$,
form a configuration on the same input with state $q$, the same input-head position and
output, and work-tape contents changed only on tapes $1$ and $3$, where they are set to
the buffer encodings of $b$ and $e$ relative to their respective old head positions;
shift work-tape head $0$ by $p$ and leave all other work-tape head positions unchanged.
\end{definition}

\begin{theorem}[Empty Snapshot on Blank Tapes]
\lean{Turing.recordStepSnapshot_empty}
\label{Turing.recordStepSnapshot_empty}
\leanok
\uses{Turing.Cfg, Turing.recordStepSnapshot}
\difficulty{4}
Let $d$ be a configuration with four work tapes and let $q$ be a state. Suppose that
every tape other than tape $0$ is blank at every integer cell. Then the snapshot formed
from $d$ using state $q$, displacement zero, and two empty buffer words is exactly the
configuration with state $q$, the same input-head position, all the same work-tape
contents and head positions, and the same output as $d$.
\end{theorem}

\begin{theorem}[Reading a Recorded Prefix]
\lean{Turing.recordStep_read_prefix}
\label{Turing.recordStep_read_prefix}
\leanok
\uses{Turing.Cfg, Turing.FinTM.bufferTape, Turing.MultiTapeTM, Turing.MultiTapeTM.runFrom, Turing.RecSym, Turing.RecSym.digitEmb, Turing.RecordAt, Turing.RecordRun, Turing.RecordStepControl, Turing.RecordStepData, Turing.binDigits, Turing.cfgRecord, Turing.cfgRecord_length, Turing.copyTM, Turing.copyTM_run_ofCfg, Turing.fieldOffset, Turing.fieldOffset_add_length_le, Turing.fieldsTape_field, Turing.joinFields, Turing.posWidth, Turing.recordRun_trans, Turing.recordStepEntry, Turing.recordStepMachine, Turing.recordStepPosFold, Turing.recordStepRows, Turing.recordStepSnapshot, Turing.recordStepSnapshot_empty, Turing.recordStepVerdict, Turing.recordStep_phase_run, Turing.recordStep_position_fold, Turing.recordStep_scan_preserves, Turing.recordStep_state_run, Turing.record_digits_value, Turing.record_pos_lt, Turing.scanFold, Turing.scanKeep, Turing.scanResult, Turing.scanTM, Turing.skipTM, Turing.skipTM_run_ofCfg}
\difficulty{6}
Let $M$ be a $k$-work-tape Turing machine on Boolean symbols with state set $Q$, let $T$
be a tag set, let $u:T\times\{\text{blank},0,1\}\to T$ be a tag update rule, and let $S$
be the state set of a two-work-tape seeking machine with designated states $s_0,s_1$.
For a Boolean input $x$, configuration $c$ of $M$ on $x$, tag $\tau\in T$, and
four-work-tape record-alphabet configuration $d$, suppose the control of $d$ is the
entry control and tape $0$ of $d$ agrees, at every offset from $-1$ through the
joined-record length, with the field encoding of $c$ and $\tau$; suppose also that all
other work tapes of $d$ are blank. If $w$ is the binary-digit encoding of the input-head
position of $c$, the record-processing machine has a run of exactly $6|w|+17$ steps from
$d$ to the snapshot whose control is at the position-skip stage, whose position-field
head is shifted by $2$, and whose two buffers both contain $w$. Before the final step
the run visits neither verdict state, and every work-tape head remains within distance
$L+1$ of its initial position, where $L$ is the length of the joined field encoding.
\end{theorem}

\begin{definition}[Record Step Read Data]
\lean{Turing.recordStepReadData}
\label{Turing.recordStepReadData}
\leanok
\uses{Turing.Cfg, Turing.RecordStepData, Turing.posWidth}
Let $x$ be a Boolean input list, $c$ a configuration with $k$ work tapes and state set
$Q$, and $s,r\in\mathbb{N}$. The constructed record stores $c$'s state, whether the
positional width of $x$ is one, and whether the shifted input-head position is zero. For
each work tape $i$, it stores the symbol under that tape's head if $i<r$ and otherwise
stores blank; it also records whether each tape index $i<r$ has its head at $-s$ or at
$s$. The tag, native-right flag, and input-symbol fields are left at their default
values.
\end{definition}

\begin{theorem}[Successor Update of Read Data]
\lean{Turing.recordStepReadData_succ}
\label{Turing.recordStepReadData_succ}
\leanok
\uses{Turing.Cfg, Turing.recordStepReadData}
\difficulty{3}
Let $x$ be a Boolean input word, let $s$ be a natural number, let $c$ be a machine
configuration on $x$, and let $i$ be a work-tape index. The read-data record at
threshold $i+1$ is obtained from the record at threshold $i$ by updating its $i$th
work-tape symbol entry to the symbol scanned by tape $i$, and its left- and right-edge
entries to indicate respectively whether that tape's head is at $-s$ or at $s$.
\end{theorem}

\begin{definition}[Record Step Read Exit]
\lean{Turing.recordStepReadExit}
\label{Turing.recordStepReadExit}
\leanok
\uses{Turing.RecordStepControl, Turing.RecordStepData}
Given stored record data with state in an optional state type, a width-one flag, and
other record fields, the resulting control state is the sweep phase for clearing the
halting case if the state is absent. Otherwise, if the width-one flag is true, it is the
sweep phase for clearing the empty case; if neither condition holds, it is the forward
scan phase at the fuel stage.
\end{definition}

\begin{theorem}[Reading Fields from a Recorded Configuration]
\lean{Turing.recordStep_read_fields}
\label{Turing.recordStep_read_fields}
\leanok
\uses{Turing.Cfg, Turing.Cfg.InWindow, Turing.FinTM.bufferTape, Turing.MultiTapeTM, Turing.MultiTapeTM.runFrom, Turing.RecSym, Turing.RecSym.digitEmb, Turing.RecordAt, Turing.RecordRun, Turing.RecordStepControl, Turing.RecordStepData, Turing.binDigits, Turing.cfgRecord, Turing.cfgRecord_length, Turing.fieldOffset, Turing.fieldOffset_add_length_le, Turing.fieldOffset_cons_succ, Turing.fieldOffset_succ, Turing.fieldsTape_field, Turing.joinFields, Turing.posWidth, Turing.recordRun_trans, Turing.recordRun_zero, Turing.recordStepFirstTape, Turing.recordStepMachine, Turing.recordStepNext, Turing.recordStepReadData, Turing.recordStepReadData_succ, Turing.recordStepReadExit, Turing.recordStepRows, Turing.recordStepSnapshot, Turing.recordStepTapeFold, Turing.recordStepVerdict, Turing.recordStep_phase_run, Turing.recordStep_scan_preserves, Turing.recordStep_tape_fold, Turing.scanFold, Turing.scanKeep, Turing.scanResult, Turing.scanResult_fold, Turing.scanTM, Turing.skipBackTM, Turing.skipBackTM_run_ofCfg, Turing.skipTM, Turing.skipTM_run_ofCfg}
\difficulty{7}
Let $M$ be a Turing machine with $k$ work tapes, Boolean work symbols, and state set
$Q$, let $T$ be a tag set, and let $u:T\times\operatorname{Option}(\{0,1\})\to T$ be a
tag update function. Let $S$ be the state set of a two-work-tape Turing machine with
designated states $s_0,s_1$. For a Boolean input $x$, a configuration $c$ of $M$ on $x$,
a tag $\tau\in T$, and $s\in\mathbb N$, assume every work-tape head of $c$ is within
distance $s$ of the origin and every work-tape cell outside that range is blank. If $d$
is a four-work-tape configuration over the record alphabet whose tape $0$, at every
offset from $-1$ through the length of the joined encoding of $c$ and $\tau$, contains
that encoding, then the record-processing machine has a run of at most $10(k+3)(L+1)$
steps, where $L$ is the length of that joined encoding. The run starts from the
specified snapshot with control at the position-skip stage and both buffers containing
the binary digits of the input-head position, and ends at the specified read-exit
snapshot, whose recorded data is read from $c$ together with the tag $\tau$; before its
final step it avoids both verdict states, and throughout it every work-tape head remains
within distance $L+1$ of its initial position.
\end{theorem}

\begin{theorem}[Reading a complete configuration record]
\lean{Turing.recordStep_read_all}
\label{Turing.recordStep_read_all}
\leanok
\uses{Turing.Cfg, Turing.Cfg.InWindow, Turing.MultiTapeTM, Turing.RecSym, Turing.RecSym.digitEmb, Turing.RecordAt, Turing.RecordRun, Turing.RecordStepControl, Turing.binDigits, Turing.cfgRecord, Turing.cfgRecord_length, Turing.fieldOffset_add_length_le, Turing.joinFields, Turing.posWidth, Turing.recordRun_trans, Turing.recordStepEntry, Turing.recordStepMachine, Turing.recordStepReadData, Turing.recordStepReadExit, Turing.recordStepSnapshot, Turing.recordStepVerdict, Turing.recordStep_read_fields, Turing.recordStep_read_prefix}
\difficulty{4}
Let $k$ be natural, let $Q,T,S$ be types with decidable equality, let $M$ be a $k$-tape
Boolean Turing machine, let $u:T\times\mathrm{Option}(\mathrm{Bool})\to T$, let $N$ be a
two-tape Boolean machine with fixed states $s_0,s_1\in S$, and let $x$ be a Boolean
input word. For a natural number $s$, a configuration $c$ of $M$, a tag $\tau\in T$, and
a four-tape record-machine configuration $d$, assume every work-tape head of $c$ lies
within distance $s$ of zero and every tape cell outside that window is blank; assume
that $d$ records the full configuration record formed from $c$ and $\tau$ on tape 0, and
that all other tapes of $d$ are blank. Then the record machine reaches the snapshot
whose control is the read-exit control associated with the data read from $c$ (with tag
$\tau$), and whose tapes 1 and 3 each hold the encoded binary input-position digits,
within $(10(k+3)+17)(L+1)$ steps, where $L$ is the length of the joined configuration
record and each work-tape head stays within distance $L+1$ of its initial position. The
run avoids both verdict states before its final time.
\end{theorem}

\begin{theorem}[Clear-Tape Snapshot and Head Bounds]
\lean{Turing.recordStep_clear_snapshot}
\label{Turing.recordStep_clear_snapshot}
\leanok
\uses{Turing.Cfg, Turing.FinTM.bufferTape, Turing.MultiTapeTM.runFrom, Turing.clearTM, Turing.clearTM_run_ofCfg, Turing.recordStepSnapshot}
\difficulty{4}
Let $d$ be a configuration of a four-work-tape machine, let $p$ be an integer, let
$b,e,w$ be finite words over its symbol alphabet, and let $j_0$ be one of the four
tapes. Suppose $j_0$ is tape $1$ or tape $3$, with $b=w$ if $j_0=1$ and $e=w$ if
$j_0=3$. Start the tape-clearing machine on the snapshot of $d$ in phase `sweep`, with
offset $p$ and buffers $b,e$. After $2|w|+2$ steps, the resulting configuration is the
snapshot in phase `done` with offset $p$, with the destination buffer emptied and the
other buffer unchanged. At every earlier step $u<2|w|+2$ the state is not `done`; at
every step $u\le 2|w|+2$, the destination head lies in the integer interval from its
position in $d$ minus $1$ through that position plus $|w|$, while every other head
remains at its initial snapshot position.
\end{theorem}

\begin{theorem}[Clear-Tape Head Displacement Bound]
\lean{Turing.recordStep_clear_heads}
\label{Turing.recordStep_clear_heads}
\leanok
\uses{Turing.Cfg, Turing.MultiTapeTM.runFrom, Turing.clearTM, Turing.recordStepSnapshot, Turing.recordStep_clear_snapshot}
\difficulty{3}
Let $d$ be a configuration of a four-work-tape machine, let $b,e,w$ be finite words over
its symbol alphabet, and let $j_0$ be one of its four tapes. Suppose $j_0$ is tape $1$
or tape $3$, with $b=w$ if $j_0=1$ and $e=w$ if $j_0=3$. Starting from the snapshot with
phase `sweep`, zero offset, and buffer words $b,e$, after any $u\le 2|w|+2$ steps of the
machine that clears tape $j_0$, every work-tape head lies within distance $|w|+1$ of its
position in $d$.
\end{theorem}

\begin{theorem}[Halted-State Record Run]
\lean{Turing.recordStep_halt_run}
\label{Turing.recordStep_halt_run}
\leanok
\uses{Turing.Cfg, Turing.Cfg.mapState, Turing.MultiTapeTM, Turing.RecSym, Turing.RecordRun, Turing.RecordStepControl, Turing.RecordStepData, Turing.recordRun_trans, Turing.recordStepMachine, Turing.recordStepReadExit, Turing.recordStepSnapshot, Turing.recordStepVerdict, Turing.recordStep_clear_heads, Turing.recordStep_clear_snapshot, Turing.recordStep_phase_run}
\difficulty{5}
Let $M$ be a deterministic machine with $k$ Boolean work tapes and state set $Q$, let
$T$ be a tag type, and let $S$ be the state type of a two-work-tape seeking machine.
Choose a tag update function, the seeking machine, and its start and finish states; take
an encoded input word $x$, a four-tape configuration $d$, a record-symbol word $w$, and
record data whose state is empty. From the snapshot in the read-exit phase with both
buffers equal to $w$, the assembled record-step machine reaches, after $4|w|+6$ steps,
the snapshot in the true-verdict phase with empty buffers. Before the final step it
reaches neither designated verdict state, and at every step through the run each
work-tape head stays within distance $|w|+1$ of its corresponding position in $d$.
\end{theorem}

\begin{theorem}[Transporting a bounded seek computation]
\lean{Turing.recordStep_seek_transport}
\label{Turing.recordStep_seek_transport}
\leanok
\uses{Turing.Cfg, Turing.Cfg.liftSymbols, Turing.FinTM.bufferTape, Turing.FinTM.bufferTape_map_decode, Turing.MultiTapeTM, Turing.MultiTapeTM.liftSymbols, Turing.MultiTapeTM.runFrom, Turing.RecSym, Turing.RecSym.digitEmb, Turing.decodeBool, Turing.embedEmitCfg, Turing.embedEmitCfg_selected_pos, Turing.embedEmitTM, Turing.embedEmitTM_frame, Turing.recordStep_transport}
\difficulty{6}
Let $Q,T,S$ be types with decidable equality, let $\mathrm{seek}$ be a two-work-tape
Boolean machine with designated states $\mathrm{start},\mathrm{finish}$, and let $C$ be
a natural number. Suppose that for every Boolean input $x$, natural number $p$, and
Boolean word $w$, with $p\le |x|+1$, any configuration $c$ on $x$ in state
$\mathrm{start}$ whose first work tape stores $w$ around its head at offsets $-1$
through $|w|$ and whose second tape is blank has a first arrival time $t\le
C(|w|+c_{\mathrm{in}}+p+1)$ at state $\mathrm{finish}$. Until that time, both work-tape
heads remain within distance $C(|w|+1)$ of their starting positions; at time $t$, the
input-head position is $p$ and both tapes, both work-head positions, and the output are
unchanged. Now let $\iota$ embed two work-tape indices into four, and let $d$ be a
four-tape configuration over record symbols on the digit-embedded input, in state
$\mathrm{start}$. Suppose every symbol stored on either embedded tape is a digit symbol,
the first embedded tape stores the digit embedding of $w$ around its head at offsets
$-1$ through $|w|$, and the second embedded tape is blank. Then the machine obtained by
embedding the symbol-lifted seek machine has a first arrival at state $\mathrm{finish}$
at some time $t\le C(|w|+d_{\mathrm{in}}+p+1)$; its work-tape heads stay within distance
$C(|w|+1)$ of their initial positions through time $t$, and at time $t$ the entire
configuration is $d$ with state changed to $\mathrm{finish}$ and input-head position
$p$.
\end{theorem}

\begin{theorem}[Mapped scan snapshot]
\lean{Turing.recordStep_map_snapshot}
\label{Turing.recordStep_map_snapshot}
\leanok
\uses{Turing.Cfg, Turing.FinTM.bufferTape, Turing.MultiTapeTM.runFrom, Turing.recordStepSnapshot, Turing.scanKeep, Turing.scanMap, Turing.scanResult_map, Turing.scanTM, Turing.scanTM_run_ofCfg}
\difficulty{6}
Let $d$ be a four-tape configuration over an alphabet $\Gamma$, let $p\in\mathbb Z$, and
let $b,e,w$ be words over $\Gamma$. For a function $f:\Gamma\to\Gamma$ and a destination
tape $j$ equal to either tape $1$ or tape $3$, assume that $b=w$ if $j=1$ and $e=w$ if
$j=3$. Starting from the snapshot of $d$ in the forward scan state, with offset $p$ and
buffers $b,e$, the scan machine that applies $f$ to symbols on tape $j$ reaches after
$2|w|+2$ steps the done-state snapshot with offset $p$, with that tape's buffer replaced
by the word obtained by applying $f$ entrywise to $w$ and the other buffer unchanged.
Before that time it is never in a done state; at every time up to and including that
time, the head on tape $j$ lies between the initial head position minus one and the
initial head position plus $|w|$, inclusive, and every other work-tape head remains at
its position in the initial snapshot.
\end{theorem}

\begin{theorem}[Empty-Input Record Run]
\lean{Turing.recordStep_empty_run}
\label{Turing.recordStep_empty_run}
\leanok
\uses{Turing.Cfg, Turing.Cfg.mapState, Turing.MultiTapeTM, Turing.RecSym, Turing.RecordRun, Turing.RecordStepControl, Turing.RecordStepData, Turing.recordStepApply, Turing.recordStepMachine, Turing.recordStepReadExit, Turing.recordStepSnapshot, Turing.recordStepVerdict, Turing.recordStep_clear_heads, Turing.recordStep_clear_snapshot, Turing.recordStep_phase_run}
\difficulty{4}
Let $M$ be a deterministic machine with $k$ Boolean work tapes and state set $Q$, let
$T$ be a tag type, and let $S$ be the state type of a two-work-tape seeking machine.
Choose a tag update function, the seeking machine, and its start and finish states; take
an encoded input word $x$, a four-tape configuration $d$, a record-symbol word $w$, and
record data whose state is nonempty and whose width flag is true. From the snapshot in
the read-exit phase with both buffers equal to $w$, the assembled record-step machine
reaches, after $2|w|+3$ steps, the snapshot in the apply phase with empty input data and
empty buffers. Before the final step it reaches neither designated verdict state, and at
every step through the run each work-tape head stays within distance $|w|+1$ of its
corresponding position in $d$.
\end{theorem}

\begin{theorem}[Rejecting an Out-of-Range Recorded Step]
\lean{Turing.recordStep_reject_run}
\label{Turing.recordStep_reject_run}
\leanok
\uses{Turing.Cfg, Turing.Cfg.mapState, Turing.MultiTapeTM, Turing.RecSym, Turing.RecordRun, Turing.RecordStepControl, Turing.RecordStepData, Turing.recordStepApply, Turing.recordStepMachine, Turing.recordStepOutside, Turing.recordStepSnapshot, Turing.recordStepVerdict, Turing.recordStep_clear_heads, Turing.recordStep_clear_snapshot, Turing.recordStep_phase_run}
\difficulty{4}
Let $M$ be a $k$-work-tape Turing machine on Boolean symbols with state set $Q$, let $T$
be a tag set, let $u:T\times\{\text{blank},0,1\}\to T$ be a tag update rule, and let $S$
be the state set of a two-work-tape seeking machine with designated states $s_0,s_1$.
Let $d$ be a four-work-tape configuration over the record alphabet, let $w$ be a finite
record word, and let $D$ be recorded step data for $M$. If the data indicate that the
simulated step moves a work-tape head outward from a recorded boundary, then the
record-processing machine has a run of exactly $2|w|+3$ steps from the snapshot with
control given by applying the recorded step to $D$, unshifted heads, and a buffer
containing $w$ on the first buffer tape, to the snapshot with rejecting verdict control,
unshifted heads, and empty buffers. Before the final step it visits neither verdict
state, and every work-tape head remains within distance $|w|+1$ of its initial position.
\end{theorem}

\begin{theorem}[Transporting a counter-saving run]
\lean{Turing.recordStep_save_transport}
\label{Turing.recordStep_save_transport}
\leanok
\uses{Turing.Cfg, Turing.Cfg.liftSymbols, Turing.CounterLoopState, Turing.FinTM.bufferTape, Turing.FinTM.bufferTape_map, Turing.FinTM.bufferTape_map_decode, Turing.MultiTapeTM.liftSymbols, Turing.MultiTapeTM.runFrom, Turing.RecSym, Turing.RecSym.digitEmb, Turing.counterLoopCfg, Turing.counterLoopTM, Turing.counterWord, Turing.counterWord_length, Turing.counter_cfg_eq, Turing.decodeBool, Turing.embedEmitCfg, Turing.embedEmitCfg_selected_pos, Turing.embedEmitCfg_selected_tape, Turing.embedEmitTM, Turing.embedEmitTM_frame, Turing.recordStepDigitBank, Turing.recordStepSave, Turing.recordStep_save_run, Turing.recordStep_transport}
\difficulty{7}
Let $Q$ and $T$ be types with decidable equality, let $x$ be a Boolean word, and let $d$
be a four-tape configuration on the input obtained by embedding each bit of $x$ as a
digit. Suppose $d$ is in the counter-decrement phase, its input-head position $p$
satisfies $p\le |x|$, and a Boolean word $w$ has binary value $2^{|w|}-1$ with $p$
strictly less than that value. Suppose tape 3 stores the digit-embedded word $w$, with
the symbol at word index $i$ stored at tape coordinate $i$ minus the position of the
tape-3 head. Then for some $t\le 5(p+1)+2|w|$, the embedded counter-saving machine
reaches the escape state with input-head position zero and tape 3 storing the
digit-embedded word obtained by applying the fixed-width decrement operation $p+1$ times
to $w$; all other configuration components are unchanged. Before time $t$ it reaches
neither the done state nor the escape state, and at every time through $t$ each
work-tape head remains within distance $|w|+1$ of its initial position. The bitwise
complements of the resulting counter word have binary value $p+1$.
\end{theorem}

\begin{theorem}[Preparation of the native input scan]
\lean{Turing.recordStep_prepare_native}
\label{Turing.recordStep_prepare_native}
\leanok
\uses{Turing.Cfg, Turing.Cfg.mapState, Turing.CounterLoopState, Turing.MultiTapeTM, Turing.MultiTapeTM.runFrom, Turing.RecSym, Turing.RecSym.digitEmb, Turing.RecordRun, Turing.RecordStepControl, Turing.RecordStepData, Turing.counterWord, Turing.recordRun_trans, Turing.recordStepMachine, Turing.recordStepProbe, Turing.recordStepReadExit, Turing.recordStepSnapshot, Turing.recordStepVerdict, Turing.recordStep_complement_value, Turing.recordStep_map_snapshot, Turing.recordStep_phase_run, Turing.recordStep_probe_run, Turing.recordStep_save_transport, Turing.scanKeep, Turing.scanMap, Turing.scanTM}
\difficulty{7}
Let $k$ be natural, let $Q,T,S$ be types with decidable equality, let $M$ be a $k$-tape
Boolean Turing machine, let $u:T\times\mathrm{Option}(\mathrm{Bool})\to T$, let $N$ be a
two-tape Boolean machine with states in $S$, and fix $s_0,s_1\in S$. For a Boolean input
word $x$, a four-tape record configuration $d$, a Boolean word $w$, and record data
$d_0$, assume that the recorded state is nonempty, the recorded width-one flag is false,
$|x|>0$, and $|x|+1<2^{|w|}$. Then there are a Boolean value $b$ and natural numbers
$q,t$ such that $b$ records whether the input head of $d$ is at the right boundary
$|x|+1$, $q=|x|$ if $b$ and otherwise equals that head position, and $q\le |x|$. The
record machine reaches the specified snapshot with its control at the seek stage in
state $s_0$, its input head at position zero, and its tape 3 holding the counter word
obtained by decrementing the all-true word of length $|w|$ exactly $q+1$ times; the run
avoids both verdict states before time $t$, keeps every work-tape head within distance
$|w|+1$ of its initial position, and satisfies $t\le 6|w|+5|x|+12$.
\end{theorem}

\begin{theorem}[Bounded seek and read behavior]
\lean{Turing.recordStep_seek_read}
\label{Turing.recordStep_seek_read}
\leanok
\uses{Turing.Cfg, Turing.Cfg.inputSymbol, Turing.Cfg.liftSymbols, Turing.Cfg.liftSymbols_inputSymbol, Turing.Cfg.mapState, Turing.FinTM.bufferTape, Turing.FinTM.bufferTape_map, Turing.MultiTapeTM, Turing.MultiTapeTM.runFrom, Turing.RecSym, Turing.RecSym.digitEmb, Turing.RecordRun, Turing.RecordStepControl, Turing.RecordStepData, Turing.decodeBool, Turing.decodeBool_apply, Turing.recordRun_trans, Turing.recordStepMachine, Turing.recordStepNext, Turing.recordStepRead, Turing.recordStepSeekBank, Turing.recordStepSnapshot, Turing.recordStepVerdict, Turing.recordStep_native_frame, Turing.recordStep_phase_run, Turing.recordStep_read_run, Turing.recordStep_seek_transport}
\difficulty{6}
Let $Q,T,S$ be types with decidable equality, let $M$ be a Boolean multi-tape machine
with $k$ work tapes, let
$\operatorname{upd}:T\times\operatorname{Option}(\mathrm{Bool})\to T$, and let
$\mathrm{seek}$ be a two-work-tape Boolean machine with designated states
$\mathrm{start},\mathrm{finish}$ and a natural-number bound constant $C$. For every
Boolean input $x$, configuration $c$ of $M$ on $x$, and Boolean word $w$, suppose that
from every two-tape configuration $a$ on input $x$ in state $\mathrm{start}$ whose first
tape stores $w$ around its head at offsets $-1$ through $|w|$ and whose second tape is
blank, there is a time $t\le C(|w|+a_{\mathrm{in}}+c_{\mathrm{in}}+1)$ at which
$\mathrm{seek}$ first reaches $\mathrm{finish}$; through time $t$ both work-tape heads
stay within distance $C(|w|+1)$ of their initial positions, and at time $t$ the
input-head position equals that of $c$, while both tapes, both work-head positions, and
the output are unchanged. Suppose also that $d$ is a four-tape configuration over record
symbols on the digit-embedded input, tape 2 is blank, and $\mathrm{saved}$ is a
record-symbol list with $\mathrm{data}$ record data. Then there is a time $t\le
C(|w|+d_{\mathrm{in}}+c_{\mathrm{in}}+1)+3$ such that the record-step machine reaches
the specified snapshot in the seek phase from the specified seek-start snapshot, does
not reach either verdict phase earlier, and keeps every work-tape head within distance
$C(|w|+1)$ of its initial position throughout.
\end{theorem}

\begin{theorem}[Bounded Binary Arithmetic Transport]
\lean{Turing.recordStep_arithmetic_transport}
\label{Turing.recordStep_arithmetic_transport}
\leanok
\uses{Turing.Cfg, Turing.Cfg.liftSymbols, Turing.FinTM.bufferTape, Turing.FinTM.bufferTape_map, Turing.FinTM.bufferTape_map_decode, Turing.FlagPhase, Turing.MultiTapeTM.liftSymbols, Turing.MultiTapeTM.runFrom, Turing.MultiTapeTM.runFrom_zero, Turing.RecSym, Turing.RecSym.digitEmb, Turing.decFixed, Turing.decodeBool, Turing.decrementTM, Turing.decrementTM_run_succ_ofCfg, Turing.embedEmitCfg, Turing.embedEmitCfg_selected_pos, Turing.embedEmitCfg_selected_tape, Turing.embedEmitTM, Turing.embedEmitTM_frame, Turing.incFixed, Turing.incrementTM, Turing.incrementTM_run_succ_ofCfg, Turing.recordStepDigitBank, Turing.recordStep_transport}
\difficulty{6}
Let $Q,T$ be types with decidable equality, let $x$ be a Boolean word, let $increase$ be
a Boolean, let $i$ be one of three digit-bank indices, and let $w,v$ be Boolean words.
Suppose that applying fixed-width increment to $w$ when $increase$ is true, and
fixed-width decrement otherwise, returns $v$, and that $v$ and $w$ have the same length.
If a four-tape configuration $d$ is in the running state and its selected digit bank
stores the encoded word $w$ relative to that tape's head position, then the
corresponding embedded arithmetic machine reaches, within $2|w|+2$ steps, the
configuration in the done-true state with that bank storing the encoding of $v$ and all
other tape contents unchanged. Before that time it never has either done state, and
throughout the run each tape head moves by at most $|w|+1$ from its initial position.
\end{theorem}

\begin{theorem}[Monotonicity of the recorded run radius]
\lean{Turing.recordRun_mono}
\label{Turing.recordRun_mono}
\leanok
\uses{Turing.Cfg, Turing.MultiTapeTM, Turing.RecordRun}
\difficulty{1}
Let a run of a $K$-work-tape machine $R$ on input $x$ start at configuration $a$, end at
configuration $b$ after $t$ steps, and avoid the states $\mathrm{yes}$ and $\mathrm{no}$
at every time $u<t$. Suppose also that, at every time $u\le t$ and on every tape $j$,
the distance between the work-head position and a specified base position is at most
$r$. For any natural number $r'\ge r$, these same run and state-avoidance conditions
hold, and all those distances are at most $r'$.
\end{theorem}

\begin{theorem}[Digit restoration during a recorded step]
\lean{Turing.recordStep_restore_digits}
\label{Turing.recordStep_restore_digits}
\leanok
\uses{Turing.Cfg, Turing.Cfg.mapState, Turing.MultiTapeTM, Turing.MultiTapeTM.runFrom, Turing.RecSym, Turing.RecSym.digitEmb, Turing.RecordRun, Turing.RecordStepControl, Turing.RecordStepData, Turing.RecordStepIndex, Turing.binDigits, Turing.counterWord, Turing.counterWord_length, Turing.recordRun_trans, Turing.recordStepDigitBank, Turing.recordStepMachine, Turing.recordStepNext, Turing.recordStepSnapshot, Turing.recordStepVerdict, Turing.recordStep_arithmetic_transport, Turing.recordStep_complement_value, Turing.recordStep_digits_succ, Turing.recordStep_map_snapshot, Turing.recordStep_phase_run, Turing.recordStep_saved_digits, Turing.scanKeep, Turing.scanMap, Turing.scanTM}
\difficulty{6}
Let $k$ be a natural number, let $Q,T,S$ be types with decidable equality, let $M$ be a
$k$-work-tape Boolean Turing machine with states in $Q$, let
$u:T\times\mathrm{Option}(\mathrm{Bool})\to T$, and let $N$ be a two-work-tape Boolean
Turing machine with states in $S$, with designated states $s_0,s_1\in S$. For a Boolean
input word $x$, a configuration $d$ of the four-tape recording machine on the input
obtained by embedding each bit of $x$ as a digit, a word $b$ over the recording
alphabet, record data $D$, and natural numbers $W,q,p$, suppose $q+1<2^W$ and $p=q+1$ if
the recorded native move is right, while $p=q$ otherwise. Then the recording machine has
a run of some length $t\le 4W+6$ from the snapshot with control in the forward
complementing phase, tape 0 head offset zero, auxiliary words $b$ and the counter word
obtained by applying the fixed-width decrement operation $q+1$ times to $W$ ones and
embedding its bits as digits, to the snapshot with control in the restoration phase at
state $s_0$, tape 0 head offset zero, and auxiliary words $b$ and the width-$W$ binary
expansion of $p$ embedded as digits. Throughout this run, neither designated verdict
state is reached before time $t$, and every work-tape head remains within distance $W+1$
of its initial position in $d$.
\end{theorem}

\begin{definition}[Bounded Record Step and Seek Specification]
\lean{Turing.RecordStepSeekSpec}
\label{Turing.RecordStepSeekSpec}
\leanok
\uses{Turing.Cfg, Turing.FinTM.bufferTape, Turing.MultiTapeTM, Turing.MultiTapeTM.runFrom}
Let $N$ be a deterministic Turing machine with two Boolean work tapes, and let $s_0,s_f$
be states of $N$. For every Boolean input list $x$, natural number $p$, and Boolean list
$w$ whose base-two digit value is $p$, assume $p\le |x|+1$. From every configuration $a$
on input $x$ that starts in state $s_0$, has the buffer encoding of $w$ on its first
work tape at every offset $z$ with $-1\le z\le |w|$, and has a blank second work tape at
every integer position, there is a time $t\le C(|w|+a.\mathrm{inputPos}+p+1)$ such that
the machine is not in state $s_f$ at any earlier time, and through time $t$ each
work-head displacement from its initial position has absolute value at most $C(|w|+1)$.
At time $t$ the machine is in state $s_f$, its input-head position is $p$, and its work
tapes, work-head positions, and output are unchanged from configuration $a$.
\end{definition}

\begin{theorem}[Restoration during a recorded transition]
\lean{Turing.recordStep_restore_native}
\label{Turing.recordStep_restore_native}
\leanok
\uses{Turing.Cfg, Turing.Cfg.mapState, Turing.FinTM.bufferTape, Turing.FinTM.bufferTape_map, Turing.MultiTapeTM, Turing.RecSym, Turing.RecSym.digitEmb, Turing.RecordRun, Turing.RecordStepControl, Turing.RecordStepData, Turing.RecordStepSeekSpec, Turing.binDigits, Turing.counterWord, Turing.recordRun_mono, Turing.recordRun_trans, Turing.recordStepApply, Turing.recordStepMachine, Turing.recordStepNext, Turing.recordStepSeekBank, Turing.recordStepSnapshot, Turing.recordStepVerdict, Turing.recordStep_clear_heads, Turing.recordStep_clear_snapshot, Turing.recordStep_phase_run, Turing.recordStep_restore_digits, Turing.recordStep_seek_transport, Turing.record_digits_value}
\difficulty{6}
Let $k$ be a natural number, let $Q,T,S$ be types with decidable equality, let $M$ be a
Boolean Turing machine with $k$ work tapes and state set $Q$, let
$u:T\times\mathrm{Option}(\mathrm{Bool})\to T$, and let $N$ be a Boolean Turing machine
with two work tapes and state set $S$. Fix states $s_0,s_1\in S$ and a natural number
$C$, and suppose that for every Boolean input $y$, natural number $p$, Boolean word $w$
with binary value $p$ and $p\le |y|+1$, and configuration $a$ of $N$ on $y$ in state
$s_0$, if tape 0 of $a$ stores $w$ at offsets $-1$ through $|w|$ from its head and tape
1 is blank everywhere, then $N$ reaches state $s_1$ at some time $v\le
C(|w|+a_{\mathrm{in}}+p+1)$, avoids state $s_1$ before time $v$, keeps both work heads
within distance $C(|w|+1)$ of their initial positions through time $v$, and at time $v$
has input-head position $p$, with its work tapes, work-head positions, and output
unchanged. Now let $x$ be a Boolean word, let $d$ be a four-work-tape configuration on
the input obtained by embedding each Boolean symbol of $x$ into the recording alphabet,
let $r$ be an input-head position, let $b$ be a word over that alphabet, let $D$ be
record data, and let $W,q$ be natural numbers. Suppose tape 2 of $d$ is blank
everywhere, $q+1<2^W$, and the input-head position of $d$ is $q+1$ if the native
input-head move recorded by $D$ is right, and $q$ otherwise. Define the transition
selected by $D$ to be the transition of $M$ from its recorded state, using its recorded
input symbol and work-tape symbols (or the stationary halting action if its state is
absent); let the next controller stage be the clearing stage if this transition attempts
to move a work head beyond a recorded tape edge, and the state-writing stage otherwise.
There is a time $t\le C(W+r+d_{\mathrm{in}}+1)+6W+10$ at which the recording machine,
started with input-head position $r$, in its forward complementing phase with data $D$,
with the word $b$ on tape 1 and the width-$W$ word obtained by applying the fixed-width
decrement operation $q+1$ times to $W$ ones on tape 3, reaches the configuration with
data given by the selected transition, that next controller stage, the original
input-head position of $d$, word $b$ on tape 1, and a blank tape 3. Before time $t$ it
reaches neither designated verdict state, and at every time up to $t$ each work-head
position stays within distance $(C+1)(W+1)$ of its position in $d$.
\end{theorem}

\begin{theorem}[Live native-step simulation]
\lean{Turing.recordStep_live_native}
\label{Turing.recordStep_live_native}
\leanok
\uses{Turing.Cfg, Turing.Cfg.inputSymbol, Turing.MultiTapeTM, Turing.RecSym, Turing.RecSym.digitEmb, Turing.RecordRun, Turing.RecordStepControl, Turing.RecordStepData, Turing.RecordStepSeekSpec, Turing.binDigits, Turing.counterWord, Turing.posWidth, Turing.recordRun_mono, Turing.recordRun_trans, Turing.recordStepApply, Turing.recordStepMachine, Turing.recordStepReadExit, Turing.recordStepSnapshot, Turing.recordStepVerdict, Turing.recordStep_empty_run, Turing.recordStep_input_blank, Turing.recordStep_prepare_native, Turing.recordStep_restore_native, Turing.recordStep_seek_read, Turing.recordStep_width_one, Turing.record_digits_value, Turing.record_pos_lt}
\difficulty{6}
Let $M$ be a deterministic machine with $k$ Boolean work tapes, let $\mathrm{upd}$
update a tag from an optional output bit, and let a two-tape Boolean machine
$\mathrm{seek}$ satisfy the following bound $C$: for every Boolean word $w$ encoding an
integer $p$ with $p\le |x|+1$, whenever its first tape stores $w$ relative to its head,
its second tape is blank, and its state is the designated start state, it reaches the
designated finish state in at most $C(|w|+a+p+1)$ steps, where $a$ is its input-head
position; before finishing it avoids that state, all work-head displacements are at most
$C(|w|+1)$, and at finishing its input-head position is $p$ while its tapes, work-head
positions, and output are unchanged. Let $x$ be a Boolean word, $c$ a configuration of
$M$ on $x$, and $d$ a four-tape configuration over the recording alphabet on the input
obtained by embedding each bit of $x$. Let $D$ be record data whose stored state is
present, whose width flag equals the assertion that the binary width of $|x|+1$ is one,
and suppose the third work tape of $d$ is blank everywhere. Then there are a bit $b$ and
a time $t\le 20(C+1)(\mathrm{width}(|x|)+|x|+1)$ such that the record-step machine,
started from the snapshot of $d$ with the read-exit state and with both buffer tapes
containing the fixed-width binary digits of $c$'s input-head position, reaches the
snapshot with the state and data prescribed by applying $M$'s transition to $D$ with
native-right flag $b$ and input symbol read from $c$, with the input-position buffer
retained and the second buffer empty. At every earlier time it avoids both verdict
states, and throughout the run every work-head displacement from its initial position is
at most $(C+1)(\mathrm{width}(|x|)+1)$.
\end{theorem}

\begin{theorem}[Observed transition matches the native step]
\lean{Turing.recordStep_observed_action}
\label{Turing.recordStep_observed_action}
\leanok
\uses{Turing.Action.apply, Turing.Cfg, Turing.Cfg.InWindow, Turing.Cfg.inputSymbol, Turing.Cfg.workTapeSymbols, Turing.MultiTapeTM, Turing.MultiTapeTM.outputSymbol, Turing.MultiTapeTM.step, Turing.recordStepAction, Turing.recordStepOutside, Turing.recordStepReadData, Turing.recordStep_window}
\difficulty{4}
Let $M$ be a Boolean-input, Boolean-work-tape machine with state set $Q$, let $s$ be a
natural number, and let $c$ be a configuration on a Boolean input $x$. Suppose every
work head of $c$ lies within distance $s$ of the origin and every tape cell at distance
greater than $s$ is blank; also suppose $c$ is not halted. For any tag $\tau$ and bit
$b$, form record data from $c$ and the window of radius $s$, then set its tag to $\tau$,
its native-right flag to $b$, and its input field to the symbol read by $c$. The action
prescribed from this data, applied to $c$, equals $M$'s one-step successor; its output
equals the symbol emitted by $M$ on that step; and the record data's outside flag is
false exactly when the successor configuration still satisfies the same window
condition.
\end{theorem}

\begin{theorem}[Recorded tape update run]
\lean{Turing.recordStep_tape_run}
\label{Turing.recordStep_tape_run}
\leanok
\uses{Turing.Action.apply, Turing.Cfg, Turing.Cfg.InWindow, Turing.Cfg.mapState, Turing.FinTM.bufferTape, Turing.MultiTapeTM, Turing.MultiTapeTM.runFrom, Turing.RecSym, Turing.RecSym.digitEmb, Turing.RecordRun, Turing.RecordStepControl, Turing.RecordStepData, Turing.ScanPhase, Turing.recordRun_trans, Turing.recordStepAction, Turing.recordStepGetMark, Turing.recordStepMachine, Turing.recordStepRows, Turing.recordStepSetMark, Turing.recordStepVerdict, Turing.recordStepWriteCell, Turing.recordStep_phase_run, Turing.recordStep_tape_word, Turing.scanKeep, Turing.scanMap, Turing.scanResult, Turing.scanResult_map, Turing.scanShift, Turing.scanTM, Turing.scanTM_run_ofCfg}
\difficulty{6}
Let $M$ be a deterministic machine with $k$ Boolean work tapes, let
$\mathrm{upd}:T\times\mathrm{Option}(\mathrm{Bool})\to T$, and let $S$ be the state set
of a two-tape Boolean machine with designated states $\mathrm{start}$ and
$\mathrm{finish}$. For an input word $x$, a configuration $c$ of $M$, a radius
$s\in\mathbb{N}$, and data $d$ and tape index $i<k$, suppose the work-tape heads of $c$
lie within distance $s$ of the origin and all tape cells outside that window are blank.
Suppose also that tape $0$ of a four-tape recording machine, relative to its head,
stores the $2s+1$ cells of tape $i$ indexed from $-s$ through $s$, each paired with a
mark indicating whether the simulated head is there. Then the recording machine has a
run of exactly $4(2s+1)+6$ steps from the configuration whose control state is the
forward writing phase for tape $i$ to the configuration whose control state skips the
newly written tape record; throughout the run it visits neither designated verdict state
before the final step, and every work-tape head stays within distance $2s+2$ of its
initial position. At the endpoint, tape $0$ stores the updated cells and head mark
obtained by applying to $c$ the action prescribed by $M$ from $d$; every other tape cell
outside tape $0$'s interval from its initial head position through that position plus
$2s$ is unchanged.
\end{theorem}

\begin{theorem}[Copyback Run Bound]
\lean{Turing.recordStep_copyback_run}
\label{Turing.recordStep_copyback_run}
\leanok
\uses{Turing.Cfg, Turing.Cfg.mapState, Turing.FinTM.bufferTape, Turing.MultiTapeTM, Turing.MultiTapeTM.runFrom, Turing.RecSym, Turing.RecordRun, Turing.RecordStepControl, Turing.RecordStepData, Turing.copyTM, Turing.copyTM_run_ofCfg, Turing.recordRun_trans, Turing.recordStepMachine, Turing.recordStepSnapshot, Turing.recordStepVerdict, Turing.recordStep_clear_heads, Turing.recordStep_clear_snapshot, Turing.recordStep_phase_run}
\difficulty{5}
Let $M$ be a deterministic machine with $k$ Boolean work tapes and state set $Q$, let
$T$ be a type of tags, let $S$ be a state type for a two-work-tape seeking machine, and
let $\mathrm{upd}:T\times\mathrm{Option}(\mathrm{Bool})\to T$ be a tag update function.
For any encoded input word $x$, four-tape configuration $d$, word $w$ over the
record-symbol alphabet, and record data, form the initial snapshot in the copyback phase
with source buffer $w$ and empty destination buffer. After $4|w|+6$ steps of the
assembled record-step machine, the configuration is the snapshot formed from $d$ with
its first work tape replaced by the buffer containing $w$, in the skip-new-position
phase with empty buffers. At every earlier step neither designated verdict state has
been reached, and throughout the run every work-tape head stays within distance $|w|+1$
of its position in $d$.
\end{theorem}

\begin{theorem}[Writing back a record field]
\lean{Turing.recordStep_write_back}
\label{Turing.recordStep_write_back}
\leanok
\uses{Turing.Cfg, Turing.Cfg.mapState, Turing.MultiTapeTM, Turing.MultiTapeTM.runFrom, Turing.RecSym, Turing.RecordAt, Turing.RecordRun, Turing.RecordStepControl, Turing.RecordStepData, Turing.fieldOffset, Turing.fieldOffset_add_length_le, Turing.fieldOffset_succ, Turing.fieldsTape_field, Turing.joinFields, Turing.recordRun_trans, Turing.recordStepMachine, Turing.recordStepNext, Turing.recordStepSnapshot, Turing.recordStepVerdict, Turing.recordStep_phase_run, Turing.skipBackTM, Turing.skipBackTM_run_ofCfg}
\difficulty{6}
Let $M$ be a deterministic machine with $k$ Boolean work tapes, let
$\mathrm{upd}:T\times\mathrm{Option}(\mathrm{Bool})\to T$, and let $\mathrm{seek}$ be a
two-tape Boolean machine with state set $S$ and designated states
$\mathrm{start},\mathrm{finish}$. Let $d$ be a four-tape configuration over the record
alphabet, and let $f_s$ be a list of fields whose length is $k+3$ and which are stored
on tape $0$ of $d$ in the sense that the record-tape contents agree with the joined
fields at every specified offset. For any recording data, there is a time $t\le
4(k+2)(L+1)$, where $L$ is the length of the joined fields, such that the recording
machine runs from the snapshot at the beginning of the backward write pass, with the
designated field offset and empty auxiliary buffers, to the snapshot in its true-verdict
state. Before the final step it visits neither the true nor false verdict state, and all
work-tape heads remain within distance $L+1$ of their initial positions.
\end{theorem}

\begin{theorem}[Recorded Input Position Update]
\lean{Turing.recordStep_position_run}
\label{Turing.recordStep_position_run}
\leanok
\uses{Turing.Action.apply, Turing.Cfg, Turing.Cfg.inputSymbol, Turing.Cfg.mapState, Turing.FinTM.bufferTape, Turing.MultiTapeTM, Turing.MultiTapeTM.liftSymbols, Turing.MultiTapeTM.runFrom, Turing.RecSym, Turing.RecSym.digitEmb, Turing.RecordRun, Turing.RecordStepControl, Turing.RecordStepData, Turing.binDigits, Turing.decFixed, Turing.decrementTM, Turing.embedEmitTM, Turing.incFixed, Turing.incrementTM, Turing.posWidth, Turing.recordRun_trans, Turing.recordStepAction, Turing.recordStepDigitBank, Turing.recordStepMachine, Turing.recordStepSnapshot, Turing.recordStepVerdict, Turing.recordStep_arithmetic_transport, Turing.recordStep_clear_heads, Turing.recordStep_clear_snapshot, Turing.recordStep_copyback_run, Turing.recordStep_phase_run, Turing.recordStep_position_digits}
\difficulty{6}
Let $k$ be a natural number, let $Q,T,S$ be types with decidable equality, let $M$ be a
$k$-work-tape Turing machine over Boolean symbols with states $Q$, let
$u:T\times\mathrm{Option}(\mathrm{Bool})\to T$, and let $N$ be a two-work-tape Turing
machine over Boolean symbols with states $S$. Fix states $s_0,s_1\in S$, a Boolean input
word $x$, a configuration $c$ of $M$ on $x$, record data $d_0$, and a four-tape
configuration $d$ for the record machine on the input obtained by encoding each symbol
of $x$ as a digit. Suppose $d_0$ records the input symbol scanned by $c$ and whether
$c$'s input-head position is zero, and suppose tape 0 of $d$ stores, relative to its
head, the buffered binary digits of the input-head position of $c$ at every integer
offset $z$ with $-1\le z\le w$, where $w$ is the bit width of $|x|+1$. Let $a$ be the
action of $M$ determined by $d_0$, and let $e$ be the initial control selected by the
input-head movement in $a$, the recorded zero flag, and the recorded input symbol: it is
the buffer-clearing stage for a zero movement, a left movement at position zero, or a
right movement with absent input symbol away from position zero; otherwise it is the
decrement or increment stage corresponding to the movement. Then the record machine,
starting with control $e$ and tape 1 containing the buffered digits of the input-head
position of $c$, reaches in some $t\le 6w+9$ steps a configuration whose control is the
skip-new-position stage and whose tape 0 contains the buffered binary digits of the
input-head position after applying $a$ to $c$, with all other tapes as specified by the
snapshot construction. At every earlier time before $t$ its control is neither verdict
state, and at every time up to and including $t$, every work-tape head remains within
distance $w+1$ of its initial position.
\end{theorem}

\begin{theorem}[Recorded Simulation of One Step]
\lean{Turing.recordStep_apply_run}
\label{Turing.recordStep_apply_run}
\leanok
\uses{Turing.Action.apply, Turing.Cfg, Turing.Cfg.InWindow, Turing.Cfg.inputSymbol, Turing.FinTM.bufferTape, Turing.MultiTapeTM, Turing.MultiTapeTM.runFrom, Turing.RecSym, Turing.RecSym.digitEmb, Turing.RecordAt, Turing.RecordRun, Turing.RecordStepControl, Turing.RecordStepData, Turing.binDigits, Turing.cfgRecord, Turing.cfgRecord_length, Turing.fieldOffset, Turing.fieldOffset_add_length_le, Turing.fieldOffset_succ, Turing.fieldsTape, Turing.fieldsTape_field, Turing.joinFields, Turing.length_joinFields_add_one, Turing.posWidth, Turing.recordRun_trans, Turing.recordRun_zero, Turing.recordStepAction, Turing.recordStepFirstTape, Turing.recordStepMachine, Turing.recordStepNext, Turing.recordStepSnapshot, Turing.recordStepVerdict, Turing.recordStep_phase_run, Turing.recordStep_position_run, Turing.recordStep_tape_run, Turing.recordStep_write_back, Turing.scanKeep, Turing.scanMap, Turing.scanResult_map, Turing.scanTM, Turing.scanTM_run_ofCfg, Turing.skipTM, Turing.skipTM_run_ofCfg}
\difficulty{7}
Let $M$ be a Boolean-input multi-tape machine with state set $Q$, let $T$ be a type of
tags, let $upd:T\times\{\text{absent},\text{false},\text{true}\}\to T$, and let $seek$
be a two-tape Boolean machine with states $S$ and designated states $start,finish$. Fix
a Boolean input word $x$, a natural-number window radius $s$, a configuration $c$ on
$x$, a tag $\tau\in T$, and a record $data$ whose input symbol agrees with the symbol
scanned by $c$, whose zero-position flag correctly records whether the input head is at
zero, and whose tag is $\tau$. Suppose $c$ lies in the window of radius $s$. Suppose
also that a four-tape configuration $d$ records $c$ and $\tau$ on tape $0$ according to
the specified record layout, while every other tape is blank. Then there is a time $t\le
40(k+3)(L+1)$, where $k$ is the number of tapes of $M$ and $L$ is the length of the
joined record, at which the recording machine reaches the configuration that has verdict
state $true$ and whose tape $0$ contains the record of the configuration obtained by
applying to $c$ the action prescribed by $M$ from $data$, with its tag updated by $upd$
using that action's output. This run avoids both verdict states before time $t$ and
keeps every work-tape head within distance $L+1$ of its initial position through time
$t$.
\end{theorem}

\begin{theorem}[Machine for Updating a Configuration Record]
\lean{Turing.exists_recordStepTM}
\label{Turing.exists_recordStepTM}
\leanok
\uses{Turing.Action.apply, Turing.Cfg, Turing.Cfg.InWindow, Turing.Cfg.inputSymbol, Turing.MultiTapeTM, Turing.MultiTapeTM.outputSymbol, Turing.MultiTapeTM.outputSymbol_of_halt, Turing.MultiTapeTM.runFrom, Turing.MultiTapeTM.step, Turing.MultiTapeTM.step_of_halt, Turing.RecSym, Turing.RecSym.digitEmb, Turing.RecordAt, Turing.RecordRun, Turing.RecordStepControl, Turing.binDigits, Turing.cfgRecord, Turing.cfgRecord_length, Turing.exists_inputSeekTM, Turing.fieldOffset_add_length_le, Turing.fieldsTape, Turing.joinFields, Turing.posWidth, Turing.recordRun_mono, Turing.recordRun_trans, Turing.recordStepAction, Turing.recordStepApply, Turing.recordStepEntry, Turing.recordStepMachine, Turing.recordStepOutside, Turing.recordStepReadData, Turing.recordStepReadExit, Turing.recordStepSnapshot, Turing.recordStepSnapshot_empty, Turing.recordStepVerdict, Turing.recordStep_apply_run, Turing.recordStep_halt_run, Turing.recordStep_live_native, Turing.recordStep_observed_action, Turing.recordStep_read_all, Turing.recordStep_reject_run}
\difficulty{6}
Let $k$ be a natural number, let $Q,T$ be finite decidable types, let $M$ be a
deterministic $k$-work-tape machine over the binary alphabet with state type $Q$, and
let $u:T\to\operatorname{Option}(\{0,1\})\to T$ satisfy $u(\tau,\mathrm{blank})=\tau$
for every $\tau\in T$. There are a natural number $K$, a finite decidable state type
$P$, a deterministic machine $R$ with $K+1$ work tapes over the record alphabet, entry
and distinct yes/no states, and a constant $C$ such that the following holds. If a
configuration $c$ has all work-tape heads and all nonblank cells within distance $s$ of
the origin, and a configuration $d$ starts in the entry state with tape $0$ encoding the
joined record of $c$ and tag $\tau$, while its other work tapes are blank, then $R$
reaches a configuration after at most
$C(|\langle\operatorname{record}(c,\tau)\rangle|+|x|+1)$ steps without visiting either
decision state earlier; through that time every work-tape head moves by at most
$C(|\langle\operatorname{record}(c,\tau)\rangle|+1)$. If the one-step successor of $c$
remains within distance $s$ on all tapes, the final state is yes and tape $0$ contains
its joined record, with tag updated by applying $u$ to $\tau$ and the output symbol
emitted by $M$ (a blank output leaves the tag unchanged); otherwise the final
configuration is the original one with state changed to no.
\end{theorem}

\begin{definition}[Record Equality Indices]
\lean{Turing.RecordEqIndex}
\label{Turing.RecordEqIndex}
\leanok
For each natural number $k$, define an index type with four kinds of indices: comparison
indices carrying $j\in\{0,\ldots,k+2\}$; forward indices carrying $j\in\{0,\ldots,k+1\}$
and $i\in\{0,1\}$; backward indices carrying a Boolean $b$, such a $j$, and such an $i$;
and verdict indices carrying a Boolean.
\end{definition}

\begin{definition}[Record Equality Phase]
\lean{Turing.recordEqPhase}
\label{Turing.recordEqPhase}
\leanok
\uses{Turing.FlagPhase, Turing.RecordEqIndex, Turing.SeekPhase}
For each natural number $k$, assign a phase type to every record-equality index: a
comparison index receives the flag phase type, while every other index receives the seek
phase type.
\end{definition}

\begin{definition}[Record Equality Backward State]
\lean{Turing.recordEqBack}
\label{Turing.recordEqBack}
\leanok
\uses{Turing.RecordEqIndex, Turing.recordEqPhase}
For each Boolean $b$ and each $j\in\{0,\ldots,k+2\}$, define a dependent control state
as follows: if $j>0$, its index is the backward index carrying $b$, $j-1$, and $0$, and
its phase is the running seek phase; if $j=0$, its index is the verdict index carrying
$b$, again with the running seek phase.
\end{definition}

\begin{definition}[Rows for record equality]
\lean{Turing.recordEqRows}
\label{Turing.recordEqRows}
\leanok
\uses{Turing.MultiTapeTM, Turing.RecordEqIndex, Turing.compareTM, Turing.recordEqPhase, Turing.skipBackTM, Turing.skipTM}
For a natural number $k$, a symbol type $\Gamma$ with decidable equality, and an
index of one of four forms—comparison with $j<k+3$, forward pass with $j<k+2$ and $i<2$,
backward pass with a Boolean flag, $j<k+2$ and $i<2$, or verdict with a Boolean
flag—define a two-work-tape Turing machine whose state type is the phase type assigned
to that index. The comparison indices use the machine that compares tapes $0$ and $1$;
forward indices use a skip machine with parameter $i$; backward indices use a
backward-skip machine with parameter $i$; and verdict indices use a skip machine with
parameter $0$.
\end{definition}

\begin{definition}[Successor map for record equality]
\lean{Turing.recordEqNext}
\label{Turing.recordEqNext}
\leanok
\uses{Turing.RecordEqIndex, Turing.recordEqBack, Turing.recordEqPhase}
For a natural number $k$ and an index of the comparison, forward, backward, or verdict
form, define a partial successor map from the phase type assigned to that index to an
index together with its corresponding phase. At a completed comparison phase, a true
flag moves to the forward pass when $j<k+2$ and otherwise to the backward pass with true
flag; a false flag moves to the backward pass with false flag. A completed forward phase
advances $i$ from $0$ to $1$, then returns to comparison at $j+1$; a completed backward
phase advances $i$ from $0$ to $1$, then moves to the backward pass at the same $j$ with
its flag. All other comparison, forward, and backward phases have no successor, and
verdict indices have no successor in any phase.
\end{definition}

\begin{theorem}[Machine for Comparing Records]
\lean{Turing.exists_recordEqTM}
\label{Turing.exists_recordEqTM}
\leanok
\uses{Turing.Cfg, Turing.FinTM.bufferTape, Turing.MultiTapeTM, Turing.MultiTapeTM.runFrom, Turing.RecSym, Turing.RecordAt, Turing.RecordEqIndex, Turing.RecordRun, Turing.compareTM_run_ofCfg, Turing.fieldOffset, Turing.fieldOffset_add_length_le, Turing.fieldOffset_succ, Turing.fieldsTape_field, Turing.joinFields, Turing.phaseHostTM, Turing.phaseHostTM_run, Turing.recordEqBack, Turing.recordEqNext, Turing.recordEqPhase, Turing.recordEqRows, Turing.recordRun_tagged, Turing.recordRun_trans, Turing.recordRun_zero, Turing.skipBackTM, Turing.skipBackTM_run_ofCfg, Turing.skipTM, Turing.skipTM_run_ofCfg}
\difficulty{7}
For every natural number $k$ and finite decidable state types $Q$ and $T$, there are a
finite decidable state type $P$, a deterministic two-work-tape machine $R$ over the
record alphabet, distinct states $\mathrm{yes},\mathrm{no}\in P$, designated entry
state, and constant $C$ such that the following holds. For any input string, lists of
fields $f,g$ each having exactly $k+3$ fields, and initial configuration $d$ in the
entry state whose two work tapes encode $f$ and $g$ at their respective head positions,
the machine reaches the original configuration with only its state changed to
$\mathrm{yes}$ if $f=g$, and to $\mathrm{no}$ otherwise, in some time $t\le C(|\langle
f\rangle|+1)$. Before time $t$ it visits neither decision state, and through time $t$
each work-tape head stays within distance $C(|\langle f\rangle|+1)$ of its starting
position; here $\langle f\rangle$ denotes the concatenation of the fields of $f$ with
blank separators, and the tape encodings agree with that concatenation at offsets from
$-1$ through its length (and likewise for $g$).
\end{theorem}

\begin{definition}[Record test indices]
\lean{Turing.RecordTestIndex}
\label{Turing.RecordTestIndex}
\leanok
For a natural number $k$ and a type $Q$, introduce an index type with a state index;
forward indices carrying an optional element of $Q$ and an index $j<k+2$; tag indices
carrying an optional element of $Q$; backward indices carrying a Boolean and an index
$j<k+2$; and verdict indices carrying a Boolean.
\end{definition}

\begin{definition}[Phases for record tests]
\lean{Turing.recordTestPhase}
\label{Turing.recordTestPhase}
\leanok
\uses{Turing.RecordTestIndex, Turing.ScanPhase, Turing.SeekPhase}
For a natural number $k$, a type $Q$, and a record-test index, assign a phase type as
follows: state indices receive the scan phases carrying an optional optional state; tag
indices receive the scan phases carrying an optional Boolean; all forward, backward, and
verdict indices receive the seek phases.
\end{definition}

\begin{definition}[Reading a recorded state]
\lean{Turing.recordReadState}
\label{Turing.recordReadState}
\leanok
\uses{Turing.RecSym}
For types $Q$ and $T$, define a partial state-reading function from record symbols of
types $Q$ and $T$ to an optional optional state: on a state-cell symbol carrying $q$, it
returns $q$ as a present result; on every other record symbol, it returns no result. The
supplied optional optional state parameter does not affect the output.
\end{definition}

\begin{definition}[Reading a recorded tag]
\lean{Turing.recordReadTag}
\label{Turing.recordReadTag}
\leanok
\uses{Turing.RecSym}
For types $Q$ and $T$, a predicate $P$ on an optional state and a tag, an optional state
$q$, and an optional Boolean input, define a partial Boolean-valued reader on record
symbols: on a tag-cell symbol carrying $\tau$, it returns the Boolean value $P(q,\tau)$
as a present result; on every other record symbol, it returns no result. The supplied
optional Boolean parameter does not affect the output.
\end{definition}

\begin{definition}[Rows for record tests]
\lean{Turing.recordTestRows}
\label{Turing.recordTestRows}
\leanok
\uses{Turing.MultiTapeTM, Turing.RecSym, Turing.RecordTestIndex, Turing.recordReadState, Turing.recordReadTag, Turing.recordTestPhase, Turing.scanFold, Turing.scanKeep, Turing.scanTM, Turing.skipBackTM, Turing.skipTM}
For a natural number $k$, types $Q$ and $T$, a predicate $P$ on an optional state and a
tag, and a record-test index, define a one-work-tape Turing machine whose state type is
the phase type assigned to that index. The state index uses a scan machine that reads
the recorded state; a tag index uses a scan machine that tests the recorded tag against
$P$ and the state carried by that index; backward indices use a backward-skip machine;
and forward and verdict indices use skip machines.
\end{definition}

\begin{definition}[Successor map for record tests]
\lean{Turing.recordTestNext}
\label{Turing.recordTestNext}
\leanok
\uses{Turing.RecordTestIndex, Turing.recordTestPhase}
For a natural number $k$, a type $Q$, and a record-test index, define a partial
successor map from its assigned phase type to a record-test index and its phase. A
completed state scan moves to the forward index carrying the scanned optional state and
$j=0$. A completed forward scan increments $j$ while $j+1<k+2$, and otherwise moves to
the tag index carrying the same state and begins its forward scan; a completed tag scan
moves to the backward index carrying the scanned Boolean and $j=k+1$. A completed
backward phase decrements $j$ while $j\ne0$, and at $j=0$ moves to the verdict index
with the same Boolean. All other state, forward, tag, and backward phases have no
successor, and verdict indices have no successor in any phase.
\end{definition}

\begin{theorem}[Machine for Testing a Configuration Record]
\lean{Turing.exists_recordTestTM}
\label{Turing.exists_recordTestTM}
\leanok
\uses{Turing.Cfg, Turing.FinTM.bufferTape, Turing.MultiTapeTM, Turing.MultiTapeTM.runFrom, Turing.RecSym, Turing.RecordAt, Turing.RecordRun, Turing.RecordTestIndex, Turing.cfgRecord, Turing.cfgRecord_length, Turing.fieldOffset, Turing.fieldOffset_add_length_le, Turing.fieldOffset_succ, Turing.fieldsTape_field, Turing.joinFields, Turing.phaseHostTM, Turing.phaseHostTM_run, Turing.recordReadState, Turing.recordReadTag, Turing.recordRun_tagged, Turing.recordRun_trans, Turing.recordRun_zero, Turing.recordTestNext, Turing.recordTestPhase, Turing.recordTestRows, Turing.scanFold, Turing.scanKeep, Turing.scanResult_fold, Turing.scanTM, Turing.scanTM_run_ofCfg, Turing.skipBackTM, Turing.skipBackTM_run_ofCfg, Turing.skipTM, Turing.skipTM_run_ofCfg}
\difficulty{7}
For finite decidable types $Q,T$ and any predicate $P$ on an optional state in $Q$ and a
tag in $T$, there are a finite decidable state type $P'$, a deterministic one-work-tape
machine $R$ over the record alphabet, entry and distinct yes/no states, and a constant
$C$ with the following property. Given any binary input $x$, natural number $s$,
configuration $c$ over $x$, tag $\tau$, and initial configuration $d$ in the entry state
whose work tape encodes the joined record of $c$ and $\tau$, the machine reaches after
at most $C(|\langle\operatorname{record}(c,\tau)\rangle|+1)$ steps the same
configuration except that its state is yes if $P$ holds of the state of $c$ and tag
$\tau$, and no otherwise. It visits neither decision state earlier, and its work-tape
head remains within distance $C(|\langle\operatorname{record}(c,\tau)\rangle|+1)$ of its
initial position throughout.
\end{theorem}

\begin{theorem}[Initial Configuration Lies in the Window]
\lean{Turing.Cfg.inWindow_init}
\label{Turing.Cfg.inWindow_init}
\leanok
\uses{Turing.Cfg, Turing.Cfg.InWindow, Turing.Cfg.init}
\difficulty{1}
Let $k$ be a nonnegative integer, let $q_0$ be a state, let $x$ be an input word, and
let $s$ be a nonnegative integer. In the initial configuration on $x$ with starting
state $q_0$, every work-tape head is at a position of absolute value at most $s$, and
every work-tape cell at an integer position $z$ with $|z|>s$ is blank.
\end{theorem}
